\documentclass[aps,prd,twocolumn,superscriptaddress,nofootinbib]{revtex4-2}

\usepackage{amsmath}
\usepackage{amssymb}
\usepackage{amsthm}
\usepackage{booktabs}
\usepackage{xcolor}
\usepackage{graphicx}
\usepackage{mathrsfs}

\newtheorem{theorem}{Theorem}[section]
\newtheorem{proposition}[theorem]{Proposition}
\newtheorem{lemma}[theorem]{Lemma}
\newtheorem{corollary}[theorem]{Corollary}
\newtheorem{conjecture}[theorem]{Conjecture}
\newtheorem{axiom}[theorem]{Axiom}
\theoremstyle{definition}
\newtheorem{definition}[theorem]{Definition}
\theoremstyle{remark}
\newtheorem{remark}[theorem]{Remark}
\newtheorem{example}[theorem]{Example}

\DeclareMathOperator{\ord}{ord}
\DeclareMathOperator{\rank}{rank}
\DeclareMathOperator{\Sel}{Sel}
\DeclareMathOperator{\Reg}{Reg}
\DeclareMathOperator{\Sha}{Sha}
\DeclareMathOperator{\Gal}{Gal}
\DeclareMathOperator{\Hom}{Hom}

\usepackage[hyperfootnotes=false]{hyperref}
\hypersetup{
    pdftitle={Die BSD-Vermutung als Positivitäts-Normalform: Eine strukturelle Übersicht},
    pdfauthor={Lukas Geiger},
    pdfsubject={Strukturelle Übersicht zur Birch-und-Swinnerton-Dyer-Vermutung über Positivitäts-Normalformen},
    pdfkeywords={Birch und Swinnerton-Dyer, Neron-Tate-Höhe, Positivität, Euler-Systeme, Gross-Zagier, Iwasawa-Theorie},
    colorlinks=true,
    linkcolor=black,
    urlcolor=blue,
    citecolor=black
}
\makeatletter
\renewcommand*{\HyperDestNameFilter}[1]{de.#1}
\makeatother

\begin{document}

% ===================================================================
% TITELSEITE
% ===================================================================

\title{\texorpdfstring{Die BSD-Vermutung als Positivitäts-Normalform:\\Eine strukturelle Übersicht}{Die BSD-Vermutung als Positivitäts-Normalform: Eine strukturelle Übersicht}}

\author{Lukas Geiger}
\affiliation{Unabhängiger Forscher, Bernau, Deutschland\\ORCID: \href{https://orcid.org/0009-0005-7296-1534}{0009-0005-7296-1534}}

\date{\today}

\begin{abstract}
Wir formulieren die Birch-und-Swinnerton-Dyer-Vermutung als \textit{Positivitäts-Normalform} um: eine strukturelle Charakterisierung, kein Beweis. Vier axiomatische Zwangsbedingungen -- endliche arithmetische Kapazität, kooperative Höhenpositivität, kontrolliertes Iwasawa-Wachstum und Euler-System-Kompositionskohärenz -- identifizieren die BSD-Führungstermformel als erwartetes ,,Sättigungsprofil'', sobald die fehlenden Brückeneingaben zusätzlich angenommen werden: Höhen-/Zykelkopplung, hinreichend starke Selmer-/Euler-System-Kontrolle und Endlichkeit von $\Sha$ in allgemeinem Rang. Für Rang~$\geq 2$ bleiben diese Brücken konjektural, und einige Axiome überlappen selbst mit BSD. Die N\'eron-Tate-Höhenpaarung liefert in dieser Lesart die kanonische Positivitätsstruktur. Wir zeigen, dass das Gross-Zagier-Theorem die Rang-1-Instantiierung dieses Frameworks ist und eine analytische Ableitung an eine nichtnegative Höhe koppelt. Für allgemeinen Rang~$r$ formulieren wir die \textit{Höhensättigungsvermutung}: $L^{(r)}(E,1)$ ist proportional zum Regulator $R_E = \det(\langle P_i, P_j \rangle) > 0$. Dies würde $\ord_{s=1} L(E,s) = \rank E(\mathbb{Q})$ erzwingen, aber erst nachdem die fehlenden höheren Gross-Zagier- und höheren Kolyvagin-/Sha-Eingaben geliefert sind.
\newpage
\end{abstract}

\keywords{Birch-und-Swinnerton-Dyer-Vermutung, N\'eron-Tate-Höhe, Positivität, Euler-Systeme, Gross-Zagier, Iwasawa-Theorie, Normalform-Charakterisierung}

\maketitle


% ===================================================================
% 1. EINFUEHRUNG
% ===================================================================
\section{Einführung}
\label{sec:intro}

\subsection{Die BSD-Vermutung}

Sei $E/\mathbb{Q}$ eine elliptische Kurve mit Führer $N$ und $L$-Funktion $L(E,s) = \sum_{n \geq 1} a_n n^{-s}$ (die für alle $E/\mathbb{Q}$ nach dem Modularitätssatz \cite{Wiles1995} existiert). Die \textit{Birch-und-Swinnerton-Dyer-Vermutung} (BSD) \cite{BSD1965} behauptet:

\begin{conjecture}[BSD -- Schwache Form]
\label{conj:bsd_weak}
$\ord_{s=1} L(E,s) = \rank E(\mathbb{Q})$.
\end{conjecture}

\begin{conjecture}[BSD -- Starke Form]
\label{conj:bsd_strong}
Der Leitkoeffizient von $L(E,s)$ bei $s = 1$ erfüllt:
\begin{equation}
\frac{L^{(r)}(E,1)}{r!} = \frac{\Omega_E \cdot R_E \cdot \prod_p c_p \cdot |\Sha(E/\mathbb{Q})|}{|E(\mathbb{Q})_{\mathrm{tors}}|^2}\,,
\label{eq:bsd_formula}
\end{equation}
wobei $r = \rank E(\mathbb{Q})$, $\Omega_E$ die reelle Periode, $R_E$ der Regulator (Determinante der N\'eron-Tate-Höhenpaarung), $c_p$ die Tamagawa-Zahlen und $\Sha(E/\mathbb{Q})$ die Tate-Shafarevich-Gruppe ist. Wir schreiben durchgängig $E_{\mathrm{tors}} := E(\mathbb{Q})_{\mathrm{tors}}$ zur Kurzform.
\end{conjecture}

Die Vermutung ist bekannt für:
\begin{itemize}
\item \textbf{Rang 0:} Wenn $L(E,1) \neq 0$, dann $\rank E(\mathbb{Q}) = 0$ und $|\Sha| < \infty$ (Kolyvagin \cite{Kolyvagin1990}, unter Verwendung von Gross-Zagier \cite{GrossZagier1986}).
\item \textbf{Rang 1:} Wenn $\ord_{s=1} L(E,s) = 1$, dann $\rank E(\mathbb{Q}) = 1$ und $|\Sha| < \infty$ (Gross-Zagier + Kolyvagin).
\item \textbf{Rang $\geq 2$:} Im Allgemeinen offen.
\end{itemize}

\subsection{Die Positivitätsperspektive}

Diese Arbeit schlägt eine strukturelle Lesart vor, in der die BSD-Vermutung, ähnlich wie Teile der Hodge- und Weil-Geschichte, über \textit{Positivität} verstanden werden kann. Die zentrale Beobachtung:

\begin{quote}
\textit{Die N\'eron-Tate-Höhenpaarung ist eine kanonische positiv-definite quadratische Form auf der Mordell-Weil-Gruppe. Sie detektiert arithmetische ,,Richtungen'' -- rationale Punkte modulo Torsion -- über ihre positiven Energiekosten. Die BSD-Vermutung behauptet, dass die $L$-Funktion genau diese Richtungen ,,sieht'', nicht mehr und nicht weniger.}
\end{quote}

Die Parallele zu bestehenden Positivitätsbeweisen ist:

\begin{center}
\begin{tabular}{lll}
\toprule
Vermutung & Positive Form & Resultat \\
\midrule
Weil & Frobenius-Eigenwerte & Deligne 1974 \\
Hodge & Hodge-Riemann-Form & Offen \\
\textbf{BSD} & \textbf{N\'eron-Tate-Höhe} & \textbf{Rang $\leq 1$} \\
\bottomrule
\end{tabular}
\end{center}

Das Gross-Zagier-Theorem \cite{GrossZagier1986} (erweitert auf alle Shimura-Kurven von Yuan-Zhang-Zhang \cite{YZZ2013}) ist das kanonische Beispiel: es koppelt die erste Ableitung $L'(E,1)$ an die Höhe $\hat{h}(P_K)$ eines Heegner-Punktes und etabliert:
\begin{equation}
L'(E,1) \neq 0 \iff \hat{h}(P_K) > 0 \iff P_K \text{ Nicht-Torsion}\,.
\end{equation}
Dies ist eine \textit{Positivitätsidentität}: ein analytisches Objekt (Ableitung von $L$) gleicht einem nichtnegativen geometrischen Objekt (Höhe), und Nicht-Verschwinden auf einer Seite erzwingt Nichttrivialität auf der anderen.

\subsection{Normalform-Theoreme}

Unter einer \textit{Positivitäts-Normalform-Charakterisierung} verstehen wir ein Resultat folgenden Typs: Gegeben eine Menge struktureller Axiome (Positivität, Endlichkeit, Kompositionskohärenz) wird eine kanonische Produktbeziehung zwischen den analytischen und geometrischen Invarianten des betrachteten Objekts ausgezeichnet. Die Weil-Vermutungen für die Zetafunktion einer Varietät über $\mathbb{F}_q$ liefern den Archetyp; im BSD-Kontext bleibt die entsprechende Charakterisierung bedingt, bis die Brückenhypothesen unabhängig etabliert sind.

\subsection{Struktur dieser Arbeit}

Wir entwickeln vier Axiome (Abschnitt~\ref{sec:axioms}), die zusammen mit den später explizit gemachten Brückenhypothesen die BSD-Führungstermformel als erwartete kanonische Produkt-Normalform auswählen. Wir zeigen, wie der Fall Rang $\leq 1$ durch bekannte Theorempakete abgedeckt und im Framework neu gelesen wird (Abschnitt~\ref{sec:rank01}), identifizieren die präzisen Obstruktionen für höheren Rang (Abschnitt~\ref{sec:higher_rank}) und diskutieren die Beziehung zur Iwasawa-Theorie und zu Euler-Systemen (Abschnitt~\ref{sec:iwasawa}).

Diese Arbeit ist als strukturelle Darstellung in sich geschlossen: Definitionen und interne Konsistenzprüfungen werden hier gegeben, während die echten Rang-$\leq 1$-Eingaben die zitierten externen Theoreme von Gross--Zagier, Kolyvagin, Kato und verwandter Iwasawa-Theorie sind. Verbindungen zum breiteren FST-Programm werden als Kontext zitiert, sind aber logisch nicht erforderlich.


% ===================================================================
% 2. DIE POSITIVITAETSSTRUKTUR
% ===================================================================
\section{Die Positivitätsstruktur von BSD}
\label{sec:positivity}

\subsection{Die N\'eron-Tate-Höhe}

Sei $E/\mathbb{Q}$ eine elliptische Kurve. Die \textit{N\'eron-Tate-(kanonische) Höhe} $\hat{h}: E(\mathbb{Q}) \to \mathbb{R}_{\geq 0}$ ist die eindeutige Funktion mit:
\begin{enumerate}
\item $\hat{h}(P) \geq 0$ für alle $P \in E(\mathbb{Q})$, mit Gleichheit genau dann wenn $P$ Torsion ist.
\item $\hat{h}(nP) = n^2 \hat{h}(P)$ (quadratische Skalierung).
\item $\hat{h}$ unterscheidet sich von der naiven (Weil-) Höhe um eine beschränkte Funktion.
\end{enumerate}

Die zugehörige \textit{N\'eron-Tate-Paarung} ist:
\begin{equation}
\langle P, Q \rangle = \frac{1}{2}\left(\hat{h}(P+Q) - \hat{h}(P) - \hat{h}(Q)\right)\,.
\label{eq:NT_pairing}
\end{equation}
Insbesondere gilt $\langle P, P \rangle = \hat{h}(P)$ für alle $P \in E(\mathbb{Q})$.

\begin{theorem}[N\'eron-Tate-Positivität]
\label{thm:NT_positive}
Die Paarung $\langle \cdot, \cdot \rangle$ ist positiv semidefinit auf $E(\mathbb{Q})$ und positiv definit auf $E(\mathbb{Q})/E(\mathbb{Q})_{\mathrm{tors}}$.
\end{theorem}

Dies ist das BSD-Analogon der Hodge-Riemann-Bilinearform: eine kanonische, nichtentartete, positiv-definite quadratische Form auf dem ,,interessanten'' Teil der arithmetischen Struktur.

\subsection{Der Regulator als Positivitätsdetektor}

Seien $P_1, \ldots, P_r$ Erzeuger von $E(\mathbb{Q})/E(\mathbb{Q})_{\mathrm{tors}}$, wobei $r = \rank E(\mathbb{Q})$. Der \textit{Regulator} ist:
\begin{equation}
R_E = \det\!\left(\langle P_i, P_j \rangle\right)_{1 \leq i,j \leq r}\,.
\label{eq:regulator}
\end{equation}

\begin{proposition}[Regulatorpositivität]
\label{prop:reg_positive}
$R_E > 0$ wann immer $r \geq 1$, und $R_E = 1$ per Konvention wenn $r = 0$.
\end{proposition}

\begin{proof}
Die Matrix $(\langle P_i, P_j \rangle)$ ist die Gram-Matrix einer positiv-definiten Form. Ihre Determinante ist strikt positiv.
\end{proof}

Der Regulator ist ein \textit{Positivitätszertifikat} für die Existenz von $r$ unabhängigen rationalen Punkten. Er kann nicht verschwinden oder negativ sein: Positiv-Definitheit von $\hat{h}$ macht $R_E > 0$ zu einer automatischen Konsequenz von $\rank E(\mathbb{Q}) = r$.

\subsection{Die BSD-Formel als Positivitätsidentität}

Die starke BSD-Formel~\eqref{eq:bsd_formula} kann gelesen werden als:
\begin{equation}
\underbrace{\frac{L^{(r)}(E,1)}{r!}}_{\text{analytisch}} = \underbrace{\Omega_E}_{> 0} \cdot \underbrace{R_E}_{> 0} \cdot \underbrace{\prod c_p}_{\geq 1} \cdot \underbrace{\frac{|\Sha|}{|E_{\mathrm{tors}}|^2}}_{\geq 0}\,.
\label{eq:bsd_positivity}
\end{equation}

Jeder Faktor auf der rechten Seite ist nichtnegativ, und $R_E > 0$ wann immer $r \geq 1$. Die Formel behauptet, dass die analytische Seite -- die a priori keinen Grund hat, nichtnegativ zu sein oder die ,,richtige'' Verschwindungsordnung zu haben -- \textit{gezwungen} ist, mit dieser positiven Struktur übereinzustimmen.

Dies ist präzise das Muster eines \textit{Positivitäts-Normalform-Theorems}: das analytische Objekt wird durch die Positivität des geometrischen/arithmetischen Objekts, dem es gleich ist, eingeschränkt.


% ===================================================================
% 3. DIE VIER AXIOME
% ===================================================================
\section{Die vier Axiome}
\label{sec:axioms}

Wir formulieren vier strukturelle Bedingungen, die jede kohärente arithmetische Theorie von $L$-Funktionen und Mordell-Weil-Gruppen erfüllen muss.

\subsection{Axiom I: Endliche Arithmetische Kapazität}

\begin{axiom}[Endliche Arithmetische Kapazität]
\label{ax:finite}
Die Tate-Shafarevich-Gruppe $\Sha(E/\mathbb{Q})$ ist endlich:
\begin{equation}
|\Sha(E/\mathbb{Q})| < \infty\,.
\end{equation}
Äquivalent in Selmer-Sprache zeigt die exakte Sequenz
\[
\begin{aligned}
0 &\longrightarrow E(\mathbb{Q})\otimes \mathbb{Q}_p/\mathbb{Z}_p \\
  &\longrightarrow \Sel_{p^\infty}(E/\mathbb{Q})
  \longrightarrow \Sha(E/\mathbb{Q})[p^\infty]
  \longrightarrow 0
\end{aligned}
\]
dass Endlichkeit von $\Sha$ den $\mathbb{Z}_p$-Korang
$\mathrm{corank}_{\mathbb{Z}_p}\Sel_{p^\infty}(E/\mathbb{Q})=\rank E(\mathbb{Q})$
und einen endlichen $p$-primären Defekt für jedes $p$ impliziert. Die
Korang-Gleichheit allein ist schwächer; der nichttriviale Inhalt von
Axiom~I ist die Endlichkeit der Tate-Shafarevich-Gruppe. (Die
$n$-Selmer-Gruppen $\Sel_n(E/\mathbb{Q})$ sind stets endlich nach einem
Standardargument über die Endlichkeit der Klassengruppe und den schwachen
Mordell-Weil-Satz.)
\end{axiom}

\textit{Status.} Endlichkeit von $\Sha$ ist bekannt für Rang $\leq 1$ (Kolyvagin \cite{Kolyvagin1990}) und für bestimmte Fälle höheren Rangs. Sie wird allgemein erwartet, bleibt aber unbewiesen.

\textit{Rolle im No-Go-Argument.} Wenn $\Sha$ unendlich wäre, könnte man sich eine Kurve mit $\rank E(\mathbb{Q}) = 0$ aber $\ord_{s=1} L(E,s) = 2$ vorstellen: der ,,fehlende Rang'' wäre in $\Sha$ verborgen. Axiom~I schließt diese Lücke.

\subsection{Axiom II: Kooperative Höhenpositivität}

\begin{axiom}[Kooperative Höhenpositivität]
\label{ax:positivity}
Die N\'eron-Tate-Höhenpaarung ist positiv-definit auf $E(\mathbb{Q})/E(\mathbb{Q})_{\mathrm{tors}}$:
\begin{equation}
\langle P, P \rangle > 0 \quad \forall\, P \in E(\mathbb{Q}) \setminus E(\mathbb{Q})_{\mathrm{tors}}\,.
\end{equation}
Außerdem ist die Höhenpaarung ,,kooperativ'': Höhen unabhängiger Punkte kombinieren quadratisch über die Gram-Determinante $R_E = \det(\langle P_i, P_j \rangle) > 0$.
\end{axiom}

\textit{Status.} Unbedingt wahr (ein Theorem, keine Vermutung).

\textit{Rolle im No-Go-Argument.} Positivität stellt sicher, dass die rechte Seite von BSD eine definite Vorzeichenstruktur hat: $R_E > 0$ ist erzwungen. Jede analytische Formel für $L^{(r)}(E,1)$ muss mit dieser Positivität kompatibel sein.

\subsection{Axiom III: Kontrolliertes Iwasawa-Wachstum}

\begin{axiom}[Kontrolliertes Iwasawa-Wachstum]
\label{ax:iwasawa}
In der zyklotomischen $\mathbb{Z}_p$-Erweiterung $\mathbb{Q}_\infty/\mathbb{Q}$ wachsen die $p$-primären Selmer-Gruppen kontrolliert und asymptotisch vorhersagbar. Schreibe $X_n = \Sel_{p^\infty}(E/\mathbb{Q}_n)^\vee$ für den Pontryagin-Dual (einen endlich erzeugten $\mathbb{Z}_p$-Modul, dessen $\mathbb{Z}_p$-Rang gleich $\rank E(\mathbb{Q}_n)$ ist):
\begin{equation}
\mathrm{ord}_p\!\bigl(|X_n^{\mathrm{tors}}|\bigr) = \mu p^n + \lambda n + \nu + O(1)\,,
\label{eq:iwasawa_growth}
\end{equation}
wobei $X_n^{\mathrm{tors}}$ den $\mathbb{Z}_p$-Torsionsuntermodul von $X_n$ bezeichnet und $\mu, \lambda \geq 0$ die Iwasawa-Invarianten sind. Die Iwasawa-Hauptvermutung liefert eine $p$-adische $L$-Funktion $L_p(E,s)$, deren charakteristisches Ideal mit dem Selmer-Modul übereinstimmt.
\end{axiom}

\textit{Status.} Die Iwasawa-Hauptvermutung für elliptische Kurven ist in wichtigen Fällen etabliert, wobei verschiedene Resultate komplementäre Aspekte abdecken:
\begin{itemize}
\item \textbf{Kato \cite{Kato2004}:} Beweist eine Teilbarkeitsrichtung ($\mathrm{char}_\Lambda(\Sel^\vee) \subseteq (\mathcal{L}_p)$, d.h.\ das charakteristische Ideal des Selmer-Duals ist im vom $p$-adischen $L$-Funktion erzeugten Ideal enthalten) für alle Primzahlen~$p$, bei denen die Modulform ordinär ist und die residuale Darstellung milde Bedingungen erfüllt. Dies gilt ohne Einschränkung an den Rang und liefert eine obere Schranke für den Selmer-Modul.
\item \textbf{Skinner-Urban \cite{SkinnerUrban2014}:} Beweisen die volle Hauptvermutung (beide Teilbarkeiten, also Gleichheit) für \textit{ordinäre} Primzahlen~$p$ (d.h.\ $p \nmid a_p$), unter technischen Bedingungen einschließlich der Irreduzibilität von $\bar{\rho}_{E,p}$ und bestimmter lokaler Hypothesen an die residuale Darstellung. Dies deckt eine große Klasse elliptischer Kurven bei ordinären Primzahlen ab, gilt aber nicht bei supersingularen Primzahlen.
\end{itemize}
Für supersingulare Primzahlen erfordert die Hauptvermutung eine andere Formulierung (Pollack, Kobayashi, Sprung) und ist Gegenstand laufender Forschung. Die volle Hauptvermutung für alle Primzahlen und alle elliptischen Kurven bleibt offen.

\textit{Rolle im No-Go-Argument.} Die kontrollierte Wachstumsformel~\eqref{eq:iwasawa_growth} (Anmerkung: ,,kontrolliert'' bezieht sich auf das asymptotisch vorhersagbare Wachstum, gesteuert durch $\mu$ und $\lambda$, nicht auf strikte Monotonie, da der $O(1)$-Fehlerterm beschränkte Schwankungen erlaubt) beschränkt das Selmer-Wachstum im $\mathbb{Z}_p$-Turm. Zusammen mit der relevanten Hauptvermutung koppelt sie die $p$-adische $L$-Funktion an den Selmer-Modul. Für Rang~$\geq 2$ bleibt dies eine $p$-adische Eingabe in das Brückenproblem, kein eigenständiger archimedischer Ranggleichheitsbeweis.

\subsection{Axiom IV: Euler-System-Kompositionskohärenz}

\begin{axiom}[Euler-System-Kompositionskohärenz]
\label{ax:euler}
In den Fällen, in denen die Brücke verfügbar ist, benötigt man ein Euler-System-Paket $\{c_m\}_{m \in \mathcal{M}}$ -- eine kohärente Familie von Kohomologieklassen, indiziert durch Moduli $m$ -- das Norm-Kompatibilitätsrelationen erfüllt:
\begin{equation}
\mathrm{cores}_{m\ell/m}(c_{m\ell}) = P_\ell(\mathrm{Frob}_\ell^{-1})\, c_m\,,
\label{eq:euler_norm}
\end{equation}
wobei $P_\ell$ der Euler-Faktor bei $\ell$ und $\mathrm{cores}$ die Korestriktion ist. Wo eine hinreichend tiefe Kolyvagin-Ableitungsmaschinerie verfügbar ist, liefert dieses Paket obere Schranken für Selmer-Gruppen:
\begin{equation}
|\Sel(E/\mathbb{Q})| \leq C \cdot |\text{(Bild des Euler-Systems)}|\,,
\end{equation}
für eine berechenbare Konstante $C$. Diese letzte Schranke liegt nicht schon in der Normkohärenz selbst; in höherem Rang ist sie genau Brückenhypothese~(H2).
\end{axiom}

\textit{Status.} Euler-Systeme existieren für viele elliptische Kurven:
\begin{itemize}
\item \textbf{Heegner-Punkte} (für Rang 1): das klassische Euler-System \cite{Kolyvagin1990}.
\item \textbf{Katos Euler-System}: aus der $K$-Theorie modularer Kurven \cite{Kato2004}.
\item \textbf{Beilinson-Flach-Elemente}: für Rankin-Selberg-Faltungen \cite{LLZ2014}.
\end{itemize}

\textit{Rolle im No-Go-Argument.} Axiom~IV hält die normkohärente Kompositionsstruktur fest, die Euler-System-Klassen besitzen müssen. Wo eine hinreichend tiefe Kolyvagin-Ableitungsmaschinerie verfügbar ist, liefert diese Struktur obere Schranken für Selmer-Gruppen. In Rang~$\leq 1$ ist dies Teil des etablierten Kato/Gross--Zagier/Kolyvagin-Bildes. In Rang~$\geq 2$ ist Normkohärenz allein keine volle Selmer-Kontrolle; die fehlende vollständige Schranke ist genau Brückenhypothese~(H2).

\begin{remark}[Wurzelzahl und Parität]
\label{rem:root_number}
Die Funktionalgleichung von $L(E,s)$ bestimmt die \textit{Wurzelzahl} $w(E) = \pm 1$, die die Parität von $\ord_{s=1} L(E,s)$ einschränkt: wenn $w(E) = +1$, ist $\ord_{s=1} L(E,s)$ gerade; wenn $w(E) = -1$, ungerade. Die \textit{Paritätsvermutung} -- dass $(-1)^{\rank E(\mathbb{Q})} = w(E)$ -- ist in vielen Fällen bewiesen (Nekov\'a\v{r} \cite{Nekovar2006}; T.~und V.~Dokchitser \cite{Dokchitser2010}). Obwohl diese Paritätsbedingung nicht zu unseren vier Axiomen gehört, liefert sie eine zusätzliche Konsistenzprüfung: jedes ,,alternative Profil'', das die Paritätsvermutung verletzt, würde allein durch die Funktionalgleichung ausgeschlossen.
\end{remark}

\begin{remark}[Unabhängigkeit der Axiome]
\label{rem:axiom-independence}
Für Rang $\leq 1$ sind die BSD-Rang- und Endlichkeitsaussagen durch unabhängige Theoreme etabliert (Gross--Zagier, Kolyvagin, Kato und verwandte Euler-System-Resultate), und das Axiomenpaket kann in den üblichen Euler-System-/Iwasawa-Kontexten instanziiert werden. Das bedeutet nicht, dass jede Formulierung von Axiom~III für jede Primzahl und jede Kurve uniform bekannt ist. Für Rang $\geq 2$ sind die Axiome~I und~IV \textit{nicht} unabhängig von BSD: Endlichkeit von $\Sha$ (Axiom~I) ist selbst eine Konsequenz der vollen BSD-Vermutung, und die Existenz eines vollständigen Kolyvagin-Systems hinreichender Tiefe (Axiom~IV) ist nicht etabliert. Das bedingte Theorem (Theorem~\ref{thm:bsd_normal}) sollte daher als \textit{strukturelle Charakterisierung} gelesen werden: es identifiziert die BSD-Produktformel als erwartete Normalform innerhalb des angegebenen Brückenpakets, anstatt einen unabhängigen Beweis von BSD zu liefern.
\end{remark}

\subsection{Heuristische Parallelen}
\label{sec:heuristic}

Die vier Axiome weisen strukturelle Parallelen zum Variationsframework von~\cite{FST-Unified2026} auf. Diese Parallelen sind \textit{rein heuristisch und motivational} -- sie deuten an, warum die Axiomstruktur natürlich ist, stellen aber keine rigorosen mathematischen Verbindungen dar und werden in keinem Beweis verwendet.

\begin{itemize}
\item \textbf{Axiom~I $\leftrightarrow$ Endliche geometrische Kapazität (Axiom~A).} Die Endlichkeit von $\Sha$ ist das arithmetische Analogon endlicher Kapazität: die ,,verborgene Masse'' ist beschränkt, was verhindert, dass unendliche Diskrepanzen in der rechten Seite der BSD-Formel absorbiert werden.

\item \textbf{Axiom~II $\leftrightarrow$ Kooperative Verstärkung (Axiom~B).} Jeder unabhängige rationale Punkt kostet positive Energie (Höhe). Die quadratische Skalierung $\hat{h}(nP) = n^2 \hat{h}(P)$ spiegelt das kooperative Wachstum in der Sättigungs-ODE wider.

\item \textbf{Axiom~III $\leftrightarrow$ Kontrolliertes Wachstum (Axiom~C).} Beim Aufsteigen im $\mathbb{Z}_p$-Turm werden Torsions- und Korangdaten der Selmer-Module durch Iwasawa-Invarianten eingeschränkt. Dies ist kontrolliertes arithmetisches Wachstum, kein direkter Ausschluss aller Rangsprünge.

\item \textbf{Axiom~IV $\leftrightarrow$ Kompositionsgesetz (Axiom~D).} Norm-kompatible Schritte im Euler-System komponieren kohärent. Die Norm-Kompatibilität~\eqref{eq:euler_norm} stellt sicher, dass das Euler-System eine global konsistente Familie ist, nicht eine Sammlung von Ad-hoc-Lokalkonstruktionen.
\end{itemize}


% ===================================================================
% 4. DIE NORMALFORM-CHARAKTERISIERUNG
% ===================================================================
\section{Die Normalform-Charakterisierung}
\label{sec:normal_form}

\subsection{Aussage}

\begin{proposition}[BSD-Normalform-Kompatibilität -- Bedingtes Template]
\label{thm:bsd_normal}
\textbf{Die Annahmen zerfallen in zwei Schichten.} Die Axiome~I--IV
formulieren die Endlichkeits-, Positivitäts-, Iwasawa-Kontroll- und
Euler-System-Kohärenzbedingungen, die jede BSD-artige Führungstermformel
respektieren muss. Der Übergang von diesen Bedingungen zum BSD-Leitterm
braucht zusätzlich Brückenhypothesen: (H1) eine Gross--Zagier-artige
Höhenkopplung, (H2) höher-rangige Kolyvagin-/Selmer-Kontrolle und (H3)
Endlichkeit von $\Sha$ in allgemeinem Rang. Für Rang~$\leq 1$ wird die
nötige Brücke durch die etablierte Gross--Zagier-, Kato- und
Kolyvagin-Maschinerie geliefert. Für Rang~$\geq 2$ gilt: Wenn Axiome~I--IV
und die Brückenhypothesen (H1)--(H3) gelten, dann ist der BSD-Leitterm der
mit diesen Bedingungen kompatible kanonische Normalformausdruck. Ohne
(H1)--(H3) ist das Resultat eine strukturelle Kompatibilitätsdiagnose, kein
Beweis von BSD oder der Ranggleichheit.

Sei $E/\mathbb{Q}$ eine elliptische Kurve, die Axiome~I--IV und die Brückenhypothesen (H1)--(H3) erfüllt. Dann gilt:
\begin{enumerate}
\item Der analytische Rang gleicht dem algebraischen Rang:
\begin{equation}
\ord_{s=1} L(E,s) = \rank E(\mathbb{Q})\,.
\end{equation}
\item Der Führungsterm hat die Normalform:
\begin{equation}
\frac{L^{(r)}(E,1)}{r!} = \Omega_E \cdot R_E \cdot \prod_p c_p \cdot \frac{|\Sha(E/\mathbb{Q})|}{|E(\mathbb{Q})_{\mathrm{tors}}|^2}\,,
\end{equation}
wobei alle Faktoren auf der rechten Seite durch die Arithmetik von $E$ bestimmt sind.
\item Kein alternatives Führungstermprofil des kanonischen BSD-Produkttyps ist mit Axiomen~I--IV plus (H1)--(H3) kompatibel.
\end{enumerate}
\end{proposition}

\begin{remark}[Bedingter Status und logische Zirkularität bei Rang~$\geq 2$]
\label{rem:conditional_status}
Proposition~\ref{thm:bsd_normal} ist eine \textit{bedingte strukturelle Charakterisierung}: sie zeigt, dass die BSD-Formel die erwartete Normalform ist, \textit{gegeben} Axiome~I--IV und die Brückenhypothesen (H1)--(H3). \textbf{Sie liefert keinen unbedingten Beweis von BSD, und für Rang~$\geq 2$ ist das Argument logisch zirkulär, wenn (H1)--(H3) als bereits verfügbar behandelt werden.} Für Rang~$\leq 1$ wird die entsprechende Brücke durch Gross--Zagier- und Kolyvagin/Kato-artige Theoreme geliefert, sodass die Schlussfolgerung eine genuine strukturelle Reformulierung bekannter Resultate ist. Für Rang~$\geq 2$ ist jedoch Axiom~I/H3 (Endlichkeit von $\Sha$) selbst eine Konsequenz der vollen BSD-Vermutung, und Axiom~IV/H2 (Existenz eines vollständigen Kolyvagin-Systems hinreichender Tiefe) ist nicht unabhängig von BSD etabliert. Die Proposition hat daher die logische Struktur ,,BSD-starke Brückeneingaben + bekannte Fakten $\Rightarrow$ BSD hat die erwartete Produktnormalform'' bei Rang~$\geq 2$, was logisch gültig, aber mathematisch trivial als Beweisstrategie ist. Ihr Wert liegt in der \textit{strukturellen Analyse}: sie identifiziert die exakt fehlenden Brückeneingaben, auch wenn sie diese nicht liefert.

Die Analogie zu den Weil-Vermutungen erfordert Qualifikation: Grothendiecks Reformulierung über \'etale Kohomologie war \textit{nicht} zirkulär -- die \'etale Kohomologie war ein unabhängig konstruiertes mathematisches Objekt, und die Reformulierung bewies Rationalität und Funktionalgleichung direkt, reduzierte das Problem auf die Riemannsche Hypothese für Varietäten. Im Gegensatz dazu sind unsere Axiome für Rang~$\geq 2$ nicht unabhängig etabliert. Die Analogie gilt \textit{im Geiste} (Identifikation des korrekten strukturellen Rahmens), aber nicht \textit{in logischer Kraft}.
\end{remark}

\subsection{Beweisstruktur}

Die Beweisstrategie hat drei Komponenten, die das Sättigungstheorem spiegeln:

\textbf{Schritt 1: Positivität erzwingt das Vorzeichen.} Durch Axiom~II gilt $R_E > 0$ wann immer $r \geq 1$. Durch Axiom~I gilt $|\Sha| < \infty$, sodass die rechte Seite der BSD-Formel eine positive reelle Zahl ist (für $r \geq 1$) oder gleich $L(E,1)/\Omega_E$ mal arithmetische Faktoren (für $r = 0$). Die analytische Seite muss mit diesem Vorzeichen übereinstimmen.

\textbf{Schritt 2: Komposition erzwingt die Ordnung.} \textit{[Status: offen für Rang $\geq 2$; bewiesen für Rang $\leq 1$.]} Für Rang~0 zeigt Katos Euler-System \cite{Kato2004}, dass $L(E,1) \neq 0$ den Rang~0 erzwingt. Für Rang~1 zeigt Kolyvagins Euler-System \cite{Kolyvagin1990}, dass aus $\ord_{s=1}L(E,s)=1$ (und einem nicht-torsionalen Heegner-Punkt) $\rank E(\mathbb{Q})=1$ und $|\Sha|<\infty$ folgen. \textbf{Für Rang~$\geq 2$ ist die allgemeine Ungleichung $\rank E(\mathbb{Q}) \leq \ord_{s=1} L(E,s)$ nicht etabliert}; sie ist eine Richtung von BSD und erfordert die unbewiesene Hypothese~(H2). Die umgekehrte Ungleichung $\ord_{s=1} L(E,s) \leq \rank E(\mathbb{Q})$ erfordert die Höhenkopplung (H1). \textit{Hinweis: Beide Richtungen dieses Schritts sind für Rang~$\geq 2$ bedingt auf unbewiesene Hypothesen (siehe Bemerkung unten).}

\begin{remark}[Heuristischer Charakter der Umkehrungleichung in Schritt~2]
Die Richtung $\ord_{s=1} L(E,s) \leq \rank E(\mathbb{Q})$ stützt sich auf die vollständige BSD-Leitterm-Formel als Annahme: sie erfordert, dass $L^{(r)}(E,1)/r!$ \textit{exakt} proportional zu $R_E \cdot |\Sha| \cdot \prod c_p / |E_{\mathrm{tors}}|^2$ ist, sodass ein Verschwinden von $L$ zu Ordnung $> r$ die Gleichung $R_E = 0$ erzwingen würde, im Widerspruch zu Axiom~II. Ohne die BSD-Formel bleibt das Argument, dass ,,der führende Term schneller als der Regulator verschwindet und damit die Positivität verletzt'', heuristisch: man brauchte einen unabhängigen Beweis dafür, dass der analytische Rang den algebraischen nicht überschreiten kann, was genau der Inhalt der (unbewiesenen) schwachen BSD-Vermutung ist. Schritt~2 ist daher als bedingt auf die BSD-Formel zu verstehen, nicht als eigenständige Deduktion aus den Axiomen~I--IV allein.
\end{remark}

\textbf{Schritt 3: Eindeutigkeit der Normalform.} Bei gegebener Ranggleichheit ist die BSD-Formel der einzige Ausdruck, der kompatibel ist mit:
\begin{itemize}
\item der Funktionalgleichung und Periodennormalisierung von $L(E,s)$ (fixiert den archimedischen Periodenfaktor $\Omega_E$, sobald das BSD-artige Leittermpaket angenommen wird),
\item dem Birch-Lemma und seinen Verallgemeinerungen (bestimmt die lokalen Faktoren $c_p$),
\item dem Regulator $R_E$ (bestimmt durch die Höhenpaarung),
\item Endlichkeit von $\Sha$ (Axiom~I).
\end{itemize}

\subsection{Eindeutigkeit des Leitterm-Profils}

\begin{proposition}[Eindeutigkeit der Leitterm-Formel -- Beweisskizze]
\label{thm:uniqueness}
Sei $E/\mathbb{Q}$ eine elliptische Kurve mit $\rank E(\mathbb{Q}) = r$, die die Axiome~I--IV und die Höhenkopplungshypothese~(H1) erfüllt. Dann ist die BSD-Leitterm-Formel~\eqref{eq:bsd_formula} der \textit{einzige} Ausdruck
\[
  \frac{L^{(r)}(E,1)}{r!} = \Phi(E)
\]
der die folgenden fünf Bedingungen gleichzeitig erfüllt:
\begin{enumerate}
  \item[\textbf{(U1)}] \textbf{Analytisch-algebraische Ranggleichheit:} $\ord_{s=1} L(E,s) = \rank E(\mathbb{Q})$.
  \item[\textbf{(U2)}] \textbf{Positivität des Leitkoeffizienten:} $\Phi(E) > 0$ wann immer $r \geq 1$ (erzwungen durch Axiom~II, da jede Kopplung an $R_E > 0$ das Vorzeichen erhalten muss).
  \item[\textbf{(U3)}] \textbf{Ganzzahligkeit des $\Sha$-Beitrags:} das Verhältnis $\Phi(E)/(\Omega_E \cdot R_E \cdot \prod_p c_p)$ hat die Form $|\Sha|/|E(\mathbb{Q})_{\mathrm{tors}}|^2$ mit $|\Sha| \in \mathbb{Z}_{>0}$ (erzwungen durch Axiom~I).
  \item[\textbf{(U4)}] \textbf{Kompatibilität mit Funktionalgleichung:} $\Phi(E)$ ist kompatibel mit der $p$-adischen Interpolationsformel $L_p(E,0) \sim L(E,1)/\Omega_E^+$ und ihren höheren Ableitungen (erzwungen durch Axiom~III, die Spezialisierung der Iwasawa-Hauptvermutung).
  \item[\textbf{(U5)}] \textbf{Isogenie-Invarianz:} $\Phi(E)$ hängt nur von der Isogenieklasse von~$E$ ab, nicht von der Wahl des Weierstrass-Modells oder der Basis von $E(\mathbb{Q})$.
\end{enumerate}
\end{proposition}

\begin{proof}[Heuristisches Argument (kein rigoroser Beweis)]
Wir argumentieren \textit{heuristisch}, dass (U1)--(U5) gemeinsam $\Phi(E)$ bestimmen. \textbf{Vorbehalt:} Dieses Argument ist eine Plausibilitätsskizze, kein rigoroser Beweis. Der rigorose Rahmen für die Eindeutigkeitsfrage ist die \textit{Bloch-Kato-Vermutung} (Tamagawa-Zahl-Vermutung) \cite{BlochKato1990,FontainePerrinRiou1994}, die Tamagawa-Zahlen für Motive definiert und die BSD-Formel als Spezialfall für das Motiv $h^1(E)(1)$ ableitet und die rationale Mehrdeutigkeit der Beilinson-Vermutungen beseitigt. Unser heuristisches Argument sollte als Motivation gelesen werden, \textit{warum} die BSD-Formel die erwartete Antwort ist; der Bloch-Kato-Rahmen liefert die korrekte mathematische Umgebung.

\textit{Schritt~1: Der Positivitätsfaktor muss $R_E$ sein.} Durch (U2) gilt $\Phi(E) > 0$ für $r \geq 1$, sodass $\Phi(E)$ eine positiv-definite Invariante von $E(\mathbb{Q})/E(\mathbb{Q})_{\mathrm{tors}}$ enthalten muss. Die einzige kanonische solche Invariante ist die Gram-Determinante $R_E = \det(\langle P_i, P_j \rangle)$ der N\'eron-Tate-Höhe, die die \textit{einzige} positiv-definite Bilinearform auf $E(\mathbb{Q})/E(\mathbb{Q})_{\mathrm{tors}}$ mit quadratischer Skalierung $\hat{h}(nP) = n^2 \hat{h}(P)$, Übereinstimmung mit der Weil-Höhe bis auf $O(1)$ und Positiv-Definitheit modulo Torsion ist (vgl.\ \cite{Silverman2009}, Theorem~VIII.9.3). \textit{Anmerkung:} Dieser Schritt zeigt nur, dass \textit{falls} $\Phi(E)$ einen positiv-definiten Faktor enthält, dieser $R_E$ sein muss; er zeigt nicht, dass $\Phi(E)$ einen solchen Faktor enthalten \textit{muss}. Letzteres erfordert die Höhenkopplungshypothese~(H1).

\textit{Schritt~2: Die archimedischen und lokalen Faktoren sind bestimmt.} (U4) (Kompatibilität mit der $p$-adischen Interpolationsformel und Funktionalgleichung) bestimmt den archimedischen Faktor $\Omega_E$ und die lokalen Faktoren $c_p$ durch das Birch-Lemma und die Spezialisierung der Iwasawa-Hauptvermutung.

\textit{Schritt~3: Der verbleibende Faktor ist durch Ganzzahligkeit eingeschränkt.} Nach Extraktion von $\Omega_E \cdot R_E \cdot \prod_p c_p$ erfordert (U3), dass das verbleibende Verhältnis $\Phi(E)/(\Omega_E \cdot R_E \cdot \prod_p c_p)$ gleich $|\Sha(E/\mathbb{Q})|/|E(\mathbb{Q})_{\mathrm{tors}}|^2$ mit $|\Sha| \in \mathbb{Z}_{>0}$ ist. \textit{Anmerkung:} Dieser Schritt setzt voraus, dass die Zerlegung $\Phi(E) = \Omega_E \cdot R_E \cdot (\text{Rest})$ gültig ist, was präzise der Inhalt der BSD-Formel ist. Die ,,Ganzzahligkeits-Einschränkung'' leitet diese Zerlegung nicht unabhängig her; vielmehr ist sie \textit{konsistent mit} der BSD-Formel, sobald diese angenommen wird.

\textit{Schritt~4: Isogenie-Invarianz schließt Ad-hoc-Konstruktionen aus.} (U5) stellt sicher, dass $\Phi(E)$ nur von der Isogenieklasse abhängt, nicht von Hilfswahlentscheidungen (Basis von $E(\mathbb{Q})$, Weierstrass-Modell). Die einzelnen Größen $\Omega_E$, $R_E$, $c_p$, $|\Sha|$, $|E(\mathbb{Q})_{\mathrm{tors}}|$ sind \textit{nicht} einzeln isogenie-invariant, aber ihre Kombination in der BSD-Formel ist es, da $L(E,s)$ nur von der Isogenieklasse abhängt (Faltings). Dies erzwingt, dass $\Phi(E)$ das einzige isogenie-invariante Produkt kanonischer arithmetischer Daten ist.

Zusammenfassend machen die fünf Bedingungen die BSD-Formel zum einzigen \textit{natürlichen Kandidaten} der Form ,,Produkt kanonischer arithmetischer Invarianten''. Dies ist eine schwächere Behauptung als ein rigoroser Eindeutigkeitsbeweis: sie zeigt, dass die BSD-Formel der einzige Ausdruck \textit{innerhalb der Klasse von Produkten bekannter kanonischer Invarianten} ist, nicht dass sie die einzige Formel ist, die (U1)--(U5) im Raum aller möglichen Funktionen von~$E$ erfüllt. Die Bloch-Kato-Vermutung \cite{BlochKato1990} liefert den rigorosen motivischen Rahmen, in dem diese Eindeutigkeit präzisiert werden kann.
\end{proof}

\subsection{Ausschluss von Alternativen}

Wir zeigen, dass ,,alternative Profile'' -- hypothetische Szenarien, in denen die BSD-Formel versagt -- die Axiome verletzen:

\begin{proposition}[Konsistenzprüfung: durch die Axiome ausgeschlossene Szenarien]
\label{prop:exclusion}
Die folgenden Szenarien werden durch die Axiome ausgeschlossen. \textit{Anmerkung:} Für Rang~$\leq 1$ geben die Ausschlüsse bekannte Resultate (Kato, Kolyvagin) in der Sprache unseres Axiomsystems wieder. Für Rang~$\geq 2$ werden nur (E3) und (E4) ausgeschlossen, und beide folgen trivial aus den Axiomannahmen. Die nichttrivialen Fälle (E1) und (E2) bleiben für Rang~$\geq 2$ bedingt.

\begin{enumerate}
\item[\textbf{(E1)}] \textbf{Analytische Nullstelle ohne geometrische Richtung:} $\ord_{s=1} L(E,s) = r > \rank E(\mathbb{Q})$. Euler-System-Schranken (Axiom~IV) etablieren die \textit{umgekehrte} Ungleichung $\rank E(\mathbb{Q}) \leq \ord_{s=1} L(E,s)$ und schließen dieses Szenario daher nicht direkt aus. Der Ausschluss von (E1) erfordert Axiom~I ($|\Sha| < \infty$) kombiniert mit der Höhenkopplungshypothese~(H1): wenn die Kopplung $L^{(r)}(E,1) = C_r \cdot R_E$ gilt (wobei $r = \rank E(\mathbb{Q})$), dann würde $\ord > r$ erzwingen, dass $L^{(r)}(E,1) = 0$ (da $L$ bis zu Ordnung $> r$ verschwindet), während $R_E > 0$ (Axiom~II) und $C_r > 0$, was einen Widerspruch ergibt. Für Rang $\leq 1$ liefert Gross-Zagier diese Kopplung explizit. \textit{Anmerkung:} Für Rang $\geq 2$ ist dieser Ausschluss bedingt durch (H1) und (H2).

\item[\textbf{(E2)}] \textbf{Geometrischer Rang ohne analytische Ableitung:} $\rank E(\mathbb{Q}) = r > \ord_{s=1} L(E,s) = s$. \textit{Unter der Annahme der Leitterm-Formel} und der Euler-System-Schranke (H2) ist dies ausgeschlossen: (H2) ergibt $r = \rank E(\mathbb{Q}) \leq \ord_{s=1} L(E,s) = s$, im direkten Widerspruch zu $r > s$. Für Rang~$\leq 1$ (wo Axiom~IV unbedingt über Katos Euler-System bewiesen ist) ist (E2) somit sofort ausgeschlossen. \textbf{Für Rang~$\geq 2$ ist der Ausschluss bedingt:} er hängt von der vollen Euler-System-Schranke (H2) ab, die ihrerseits die Existenz eines Kolyvagin-Systems hinreichender Tiefe erfordert -- ein großes offenes Problem. Ohne (H2) ist die Ungleichung $\rank \leq \ord$ nicht etabliert und es ergibt sich kein Widerspruch; (E2) \textit{kann derzeit für Rang~$\geq 2$ nicht ausgeschlossen werden}. Ergänzend liefert, unter der Annahme dass die Leitterm-Formel gilt, die Höhenkopplung (H1) eine kompatible Perspektive: sie behauptet, dass $L^{(r)}(E,1)/r!$ gleich $C_r \cdot R_E$ als \textit{Leitterm-Identität} ist und die $r$-te Ableitung an den Regulator koppelt genau dann, wenn $\ord_{s=1} L(E,s) = r$; wenn $\ord = s < r$, ist die Voraussetzung der Kopplungsformel nicht erfüllt.

\item[\textbf{(E3)}] \textbf{Unendliches $\Sha$ absorbiert Diskrepanz:} $\ord_{s=1} L(E,s) = r$ aber $|\Sha| = \infty$, was eine andere rechte Seite erlaubt. Ausgeschlossen durch Axiom~I.

\item[\textbf{(E4)}] \textbf{Oszillierender Rang in $\mathbb{Z}_p$-Türmen:} Rang springt nichtmonoton durch zyklotomische Schichten. Ausgeschlossen durch Axiom~III.
\end{enumerate}
\end{proposition}


% ===================================================================
% 5. RANG 0 UND 1: DIE BEWIESENEN FAELLE
% ===================================================================
\section{\texorpdfstring{Rang $\leq 1$: Positivität in Aktion}{Rang <= 1: Positivität in Aktion}}
\label{sec:rank01}

\subsection{Rang 0: Katos Ausschluss}

\begin{proposition}[Rang-0-Sättigung]
\label{prop:rank0}
Sei $E/\mathbb{Q}$ mit $L(E,1) \neq 0$. Dann sind alle vier Axiome verifiziert, $\rank E(\mathbb{Q}) = 0$, $|\Sha| < \infty$, und die BSD-Formel gilt.
\end{proposition}

\begin{proof}[Beweisskizze]
\textit{Axiom II angewandt:} Wenn $\rank E(\mathbb{Q}) \geq 1$, existiert ein Nicht-Torsionspunkt $P$ mit $\hat{h}(P) > 0$. Die direkte Implikation $L(E,1) \neq 0 \Rightarrow \rank = 0$ wird jedoch nicht allein durch Axiom~II etabliert; sie folgt aus der Euler-System-Maschinerie (Axiom~IV) unten.

\textit{Axiom IV angewandt:} Katos Euler-System \cite{Kato2004}, konstruiert aus $K$-Theorie-Klassen auf Modulkurven, liefert die relevante obere Schranke für Rang~0. Genauer zeigt Kato für jede Primzahl~$p$, bei der die residuale Darstellung $\bar{\rho}_{E,p}$ surjektiv ist (was für alle bis auf endlich viele~$p$ gilt), dass wenn $L(E,1) \neq 0$, das Bild seiner Euler-System-Klasse in $H^1(\mathbb{Q}, T_p(E))$ nichttrivial ist, was den $\mathbb{Z}_p$-Korang der $p^\infty$-Selmer-Gruppe auf Null zwingt:
\begin{equation}
\mathrm{corank}_{\mathbb{Z}_p}\, \Sel_{p^\infty}(E/\mathbb{Q}) = 0\,,
\end{equation}
also $\rank E(\mathbb{Q}) = 0$. Kombiniert mit der Selmer-Schranke ergibt sich auch $|\Sha(E/\mathbb{Q})| < \infty$. (Anmerkung: Für Rang~0 ist das Heegner-Punkt-Euler-System nicht direkt anwendbar, da der Heegner-Punkt Torsion ist wenn $L(E,1) \neq 0$; Katos Euler-System ist das geeignete Werkzeug.)

Der Rang-0-Fall ist im Positivitäts-Framework vollständig ,,gesättigt'': alle vier Axiome sind verifiziert (Axiom~III bei allen Primzahlen guter ordinärer Reduktion mit surjektiver residualer Darstellung; die verbleibenden Primzahlen werden durch komplementäre Resultate abgedeckt oder sind endlich an der Zahl), und das No-Go-Argument schließt $\rank \geq 1$ aus.
\end{proof}

\subsection{Rang 1: Gross-Zagier als Positivitätsidentität}

Das Gross-Zagier-Theorem \cite{GrossZagier1986} besagt: für einen geeigneten imaginär-quadratischen Körper $K$, der die Heegner-Hypothese erfüllt (jede Primzahl, die den Führer $N$ teilt, zerfällt in $K$ -- ein solches $K$ existiert stets nach dem Chinesischen Restsatz und quadratischer Reziprozität), gilt
\begin{equation}
L'(E,1) = \frac{8\pi^2 \|f\|^2}{\sqrt{|D_K|}} \cdot \hat{h}(P_K)\,,
\label{eq:gross_zagier}
\end{equation}
wobei $f$ die Neuform ist, die $E$ zugeordnet ist, $\|f\|^2 = \int_{\Gamma_0(N) \backslash \mathcal{H}} |f(z)|^2 \, \frac{dx\,dy}{y^2}$ die Petersson-Norm (die $L^2$-Norm von $f$ bezüglich des hyperbolischen Masses auf der Modulkurve) bezeichnet, $D_K$ die Diskriminante und $P_K \in E(K)$ der Heegner-Punkt. (Die Formel gilt bis auf einen expliziten rationalen Faktor ungleich Null, der von $K$ abhängt; für $|D_K| > 4$ ist dieser Faktor gleich~$1$. Die Beziehung zwischen der Petersson-Norm $\|f\|^2$ und der N\'eron-Periode $\Omega_E$ in der BSD-Formel involviert die Manin-Konstante~$c_M$; für die optimale Kurve in jeder Isogenieklasse ist $c_M = 1$ rechnerisch verifiziert für alle Führer bis $500\,000$ in der Cremona-Datenbank \cite{LMFDB}. Cremonas \texttt{ecdata} liefert vollständige Tabellen bis Führer $N = 500\,000$; siehe \url{https://johncremona.github.io/ecdata/}. Diese Verifikation ist unbedingt für $N \leq 130\,000$ und bedingt auf Optimalität der ersten Kurve für $N > 400\,000$. Allgemeiner bewies \v{C}esnavi\v{c}ius~\cite{Cesnavicius2018} $c_M = 1$ für alle semistabilen elliptischen Kurven über~$\mathbb{Q}$.)

\begin{proposition}[Gross-Zagier als Positivitätsidentität]
\label{prop:gz_positivity}
Die Gross-Zagier-Formel~\eqref{eq:gross_zagier} ist eine Positivitätsidentität:
\begin{itemize}
\item Die linke Seite $L'(E,1)$ hat keine a priori Vorzeicheneinschränkung.
\item Die rechte Seite ist ein Produkt aus $\hat{h}(P_K) \geq 0$ (nach Axiom~II) und positiven Konstanten.
\item Nicht-Verschwinden: $L'(E,1) \neq 0 \iff \hat{h}(P_K) > 0 \iff P_K$ ist Nicht-Torsion.
\end{itemize}
\end{proposition}

Dies ist das BSD-Analogon der Hodge-Riemann-Bilinearrelation: das analytische Objekt gleicht einer geometrisch positiven Größe, und die Positivität erzwingt den korrekten Rang.

Kolyvagins Euler-System vervollständigt dann das Argument: es zeigt $\rank E(\mathbb{Q}) = 1$ (der Heegner-Punkt erzeugt eine endlichindizierte Untergruppe) und $|\Sha| < \infty$.

\begin{remark}
Der Rang-1-Beweis folgt exakt dem Muster eines No-Go-Theorems:
\begin{enumerate}
\item Nehme an, $\ord_{s=1} L(E,s) = 1$.
\item Gross-Zagier konvertiert das analytische Datum in ein positives geometrisches Datum ($\hat{h}(P_K) > 0$).
\item Kolyvagin schließt alternative Erklärungen aus ($\rank > 1$ oder $|\Sha| = \infty$).
\item Schlussfolgerung: $\rank E(\mathbb{Q}) = 1$, und die BSD-Formel gilt.
\end{enumerate}
Keine explizite Konstruktion rationaler Punkte wird benötigt jenseits des Heegner-Punktes (der als ,,Keim'' für das Euler-System dient, nicht als konstruktiver Beweis des Rangs).
\end{remark}

\subsection{Numerische Illustration}

\begin{example}[Kurve 11a -- Rang 0]
\label{ex:11a}
Die elliptische Kurve $E: y^2 + y = x^3 - x^2 - 10x - 20$ (Cremona-Label 11a1, Führer $N = 11$; \cite{LMFDB}) ist die Kurve mit dem kleinsten Führer. Sie hat $\rank E(\mathbb{Q}) = 0$, $E(\mathbb{Q})_{\mathrm{tors}} \cong \mathbb{Z}/5\mathbb{Z}$ und $|\Sha| = 1$ (bewiesen). Mit $R_E = 1$ (per Konvention für Rang~0), $\Omega_E = 1{,}2692\ldots$, $c_{11} = 5$ (Kodaira-Typ $I_5$, split multiplikative Reduktion) und $|E(\mathbb{Q})_{\mathrm{tors}}| = 5$:
\[
\begin{aligned}
  L(E,1)
  &= \Omega_E \cdot R_E \cdot c_{11}
     \cdot \frac{|\Sha|}{|E_{\mathrm{tors}}|^2} \\
  &= 1{,}2692 \cdot 1 \cdot 5 \cdot \frac{1}{25}
   = 0{,}2538\ldots
\end{aligned}
\]
Dies stimmt mit dem numerischen Wert $L(E,1) = 0{,}2538\ldots$ auf alle verfügbaren Stellen überein. Alle vier Axiome sind verifiziert: $|\Sha| = 1$ (I), Höhenpositivität vakuumartig erfüllt für $r = 0$ (II), Iwasawa-Hauptvermutung bewiesen bei allen Primzahlen $p \neq 11$ (III), Katos Euler-System liefert die Selmer-Schranke (IV).

Die Positivitätsinterpretation: Rang~0 ist der ,,Vakuumzustand'' des Positivitäts-Frameworks -- keine Höhenrichtung existiert, $R_E = 1$ per Konvention, und die BSD-Formel reduziert sich auf ein Verhältnis archimedischer und lokaler Daten. Die Positivitätsstruktur ist trivial gesättigt.
\end{example}

\begin{example}[Kurve 37a -- Rang 1]
\label{ex:37a}
Die elliptische Kurve $E: y^2 + y = x^3 - x$ (Cremona-Label 37a1, Führer $N = 37$) hat $\rank E(\mathbb{Q}) = 1$, erzeugt von $P = (0,0)$ mit $\hat{h}(P) = 0{,}0511\ldots$ Die BSD-Daten: $\Omega_E = 5{,}9869\ldots$, $R_E = 0{,}0511\ldots$, $c_{37} = 1$, $|E(\mathbb{Q})_{\mathrm{tors}}| = 1$, $|\Sha| = 1$ (bewiesen). Die BSD-Formel sagt vorher $L'(E,1) = \Omega_E \cdot R_E \cdot 1 \cdot 1/1 = 0{,}3059\ldots$, was mit der numerischen Berechnung $L'(E,1) = 0{,}3059\ldots$ auf alle verfügbaren Stellen übereinstimmt. Alle vier Axiome sind verifiziert: $|\Sha| = 1 < \infty$ (I), $\hat{h}(P) > 0$ (II), Iwasawa-Hauptvermutung bewiesen für $p \neq 37$ (III), Heegner-Punkt-Euler-System (IV).
\end{example}

\begin{example}[Kurve 389a -- Rang 2]
\label{ex:389a}
Die elliptische Kurve $E: y^2 + y = x^3 + x^2 - 2x$ (Cremona-Label 389a1) hat $\rank E(\mathbb{Q}) = 2$, erzeugt von $P_1 = (-1,1)$ und $P_2 = (0,0)$. Die Höhenmatrix ist:
\[
\left(\langle P_i, P_j \rangle\right) = \begin{pmatrix} 0{,}6868 & 0{,}2685 \\ 0{,}2685 & 0{,}3270 \end{pmatrix}, \quad R_E = \det = 0{,}1525\ldots
\]
Die BSD-Formel sagt vorher $L''(E,1)/2 = \Omega_E \cdot R_E \cdot \prod c_p \cdot |\Sha|/|E_{\mathrm{tors}}|^2$. Mit $\Omega_E = 4{,}9804\ldots$, $c_{389} = 1$, $|E_{\mathrm{tors}}| = 1$ und dem Referenzwert $|\Sha| = 1$ ergibt sich $L''(E,1)/2 = 0{,}7593\ldots$, übereinstimmend mit der numerischen Berechnung. Dies ist ein Rang-2-Konsistenzcheck, keine Brückenevidenz: Axiome~I und~IV sind in diesem Bereich nicht rigoros bewiesen, und die Eingabe verwendet bekannte BSD-Daten.
\end{example}


% ===================================================================
% 6. HOEHERER RANG
% ===================================================================
\section{Höherer Rang: Die offene Grenze}
\label{sec:higher_rank}

\subsection{Die Höhensättigungsvermutung}

Für allgemeinen Rang $r$ formulieren wir:

\begin{conjecture}[Höhensättigung]
\label{conj:height_saturation}
Sei $E/\mathbb{Q}$ mit $\rank E(\mathbb{Q}) = r$. Seien $P_1, \ldots, P_r$ eine Basis von $E(\mathbb{Q})/E(\mathbb{Q})_{\mathrm{tors}}$. Dann gilt:
\begin{equation}
L^{(r)}(E,1) \sim R_E = \det\!\left(\langle P_i, P_j \rangle\right) > 0\,,
\label{eq:height_saturation}
\end{equation}
in dem Sinne, dass das Verhältnis $L^{(r)}(E,1)/(r! \cdot \Omega_E \cdot R_E \cdot \prod c_p)$ gleich $|\Sha|/|E_{\mathrm{tors}}|^2$ ist.
\end{conjecture}

\begin{remark}[Beziehung zur klassischen BSD-Vermutung]
Vermutung~\ref{conj:height_saturation} ist \textit{äquivalent} zur starken BSD-Vermutung (Vermutung~\ref{conj:bsd_strong}): die ,,Höhensättigungs''-Formulierung formuliert die BSD-Leitterm-Formel lediglich in der Sprache des Positivitäts-Frameworks dieser Arbeit um. Der Name ,,Höhensättigung'' betont die Rolle des positiv-definiten Regulators $R_E > 0$ als dominantem geometrischen Faktor, aber der mathematische Inhalt ist identisch zur klassischen Formulierung. Wir behalten diese Umformulierung bei, weil sie die Positivitätsstruktur hervorhebt, nicht weil sie neuen mathematischen Inhalt liefert.
\end{remark}

Dies ist das ,,höhere Gross-Zagier'' in seiner optimistischsten Form: die $r$-te Ableitung von $L$ ist proportional zum $r \times r$-Regulator -- der Determinante einer positiv-definiten Matrix.

\begin{remark}[Cauchy-Koeffizientenformel]
\label{conj:height-saturation-identity}
Die BSD-Leitterm-Formel kann mit der Cauchy-Koeffizientenformel ausgedrückt werden, die die Rolle des Taylor-Koeffizienten explizit macht. Dies ist eine komplexanalytische Standardidentität, keine unabhängige Vermutung. Für eine elliptische Kurve $E/\mathbb{Q}$ vom analytischen Rang $r = \mathrm{ord}_{s=1} L(E,s)$:
\[
\begin{aligned}
\frac{L^{(r)}(E, 1)}{r!}
&= \frac{1}{2\pi i}
   \oint_{|s-1|=\varepsilon}
   \frac{L(E, s)}{(s-1)^{r+1}} \, ds \\
&= \Omega_E \cdot \det(\langle P_i, P_j \rangle_{\mathrm{NT}})
   \cdot \frac{\#\mathrm{Sha}(E) \cdot \prod c_p}
   {\#E(\mathbb{Q})_{\mathrm{tors}}^2}.
\end{aligned}
\]
wobei $P_1, \ldots, P_r$ Erzeuger von $E(\mathbb{Q}) / E(\mathbb{Q})_{\mathrm{tors}}$ sind und $\varepsilon>0$ klein genug gewählt wird, dass der Kreis keine Singularität des lokalen Ausdrucks außer dem Entwicklungspunkt einschließt. Das ungewichtete Integral $\oint L(E,s)\,ds$ wäre null, da $L(E,s)$ bei $s=1$ holomorph ist.
\end{remark}

\begin{remark}[Normierung der Koeffizientenextraktion]\label{rem:disc-normalisation}
Bemerkung~\ref{conj:height-saturation-identity} verwendet die Cauchy-Normierung: gilt nahe $s=1$ die Entwicklung $L(E,s)=\sum_{k\geq r} a_k(s-1)^k$, dann ist
\[
\frac{1}{2\pi i}\oint_{|s-1|=\varepsilon}
  \frac{L(E,s)}{(s-1)^{r+1}}\,ds
 = a_r = \frac{L^{(r)}(E,1)}{r!}.
\]
Äquivalent kann man den Fourier-Koeffizienten
$\varepsilon^{-r}(2\pi)^{-1}\int_0^{2\pi} L(E,1+\varepsilon e^{i\theta})e^{-ir\theta}\,d\theta$ verwenden.
Nach BSD ist dieser Koeffizient gleich $\Omega_E \cdot R_E \cdot \prod c_p \cdot |\mathrm{Sha}| / |E_{\mathrm{tors}}|^2$.
\end{remark}

\begin{remark}[Höhensättigung als Selektionsprinzip]
\label{rem:height-saturation}
Die Vermutung reformuliert die BSD-Leitterm-Formel als eine Koeffizientenidentität statt einer punktweisen Auswertung: das Cauchy-Funktional extrahiert den führenden Taylor-Koeffizienten, den BSD mit einem Produkt arithmetischer Invarianten gleichsetzt, dominiert vom positiv-definiten Regulator $R_E$. Diese Perspektive legt nahe, die BSD-Formel als \emph{Selektionsprinzip} zu betrachten: unter allen Führungstermstrukturen \textit{der Form ,,Produkt kanonischer arithmetischer Invarianten''} (siehe das heuristische Argument von Proposition~\ref{thm:uniqueness}) macht die Positivität der N\'eron-Tate-Höhenpaarung die BSD-Formel zum einzigen natürlichen Kandidaten, wobei $\Sha$ die verbleibende arithmetische Korrektur liefert. Die rigorose Formulierung dieser Eindeutigkeit wird durch die Bloch-Kato-Vermutung \cite{BlochKato1990} geliefert.
\end{remark}

\begin{remark}[Numerische Diagnostik via LMFDB]\label{rem:lmfdb-saturation}
LMFDB- und Cremona-Daten erlauben Konsistenzdiagnostik, aber keinen Ersatz für die fehlende algebraische Kontrolle. Für Kurven mit verfügbaren analytischen Leittermen und Regulatordaten kann man
\[
Q_E =
\frac{L^{(r)}(E,1)\,\#E(\mathbb{Q})_{\mathrm{tors}}^2}
     {r!\,\Omega_E\,R_E\,\prod c_p}
\]
bilden. BSD sagt voraus, dass $Q_E=\#\mathrm{Sha}(E)$ ist, also bei endlicher $\Sha$ eine nichtnegative Quadratzahlordnung. Solche Rechnungen testen Normierung, lokale Faktoren, Regulatordaten und stichprobenweise Konsistenz mit BSD. Sie beweisen weder (H1), noch (H2), noch (H3) und müssen strikt von jedem Beweisclaim getrennt bleiben. Wenn die Eingabe eine bekannte Mordell--Weil-Basis, bekannten Rang oder Regulatordaten verwendet, ist die Rechnung referenznormalisiert und muss als Advice statt als Brückenevidenz markiert werden. Smiths Verteilungsresultate für $2^\infty$-Selmer-Gruppen \cite{Smith2022} liefern einen probabilistischen Hintergrund für rangstratifizierte Stichproben.
\end{remark}

\begin{remark}[Rang-2-Höhensättigung]
\label{rem:rank2-saturation}
Für 17 elliptische Kurven vom Rang~2 aus der Cremona-Datenbank (Führer $389 \leq N \leq 5077$) erfüllt der Regulator $R > 0$ in allen Fällen ($R_{\min} = 0{,}095$, $R_{\max} = 1{,}700$). Log-Log-Regression ergibt $R \sim N^{0{,}099}$, konsistent mit sublinearem Wachstum. \textbf{Vorbehalt:} Die Stichprobengröße ($n = 17$) ist zu klein für statistisch signifikante Schlussfolgerungen; diese Resultate sollten als vorläufige Konsistenzprüfungen betrachtet werden, nicht als statistische Bestätigung eines Skalierungsgesetzes. Die Positivität $R > 0$ ist a priori durch die Positiv-Definitheit der N\'eron-Tate-Höhe garantiert (Theorem~\ref{thm:NT_positive}). Skript: \texttt{compute\_rank2\_lmfdb.py}.
\end{remark}

\subsection{Was würde ein Beweis erfordern?}

Ein Beweis der Vermutung~\ref{conj:height_saturation} über das Positivitäts-Framework erfordert drei Hypothesen, die wir zur Referenzierung im gesamten Paper kennzeichnen:

\textbf{(H1) Höhere Gross-Zagier-Formel.} Eine explizite Identität:
\begin{equation}
L^{(r)}(E,1) = C_r \cdot \det\!\left(\langle P_i, P_j \rangle\right)\,,
\label{eq:higher_gz}
\end{equation}
wobei $P_1, \ldots, P_r$ kanonische arithmetische Punkte (,,höhere Heegner-Punkte'') und $C_r > 0$ eine explizite Konstante sind.

\textit{Status:} Für $r = 1$ ist dies Gross-Zagier. Für $r = 2$ existieren Teilresultate über Diagonalrestriktion $p$-adischer Familien \cite{DarmonRotger2017}. Kim \cite{Kim2024} hat eine höhere Gross-Zagier-Formel erhalten, die die Verbindung zwischen Heegner-Punkt-Kolyvagin-Systemen und Selmer-Gruppen-Struktur verfeinert. Für allgemeines $r$ ist dies ein großes offenes Problem. Jüngere Arbeiten zu Euler-Systemen für höherrangige Motive \cite{LSZ2020} liefern Fortschritt, aber keine vollständige Formel.

\textbf{(H2) Höheres Kolyvagin-System.} Eine Euler-System-Maschinerie, die $\Sel(E/\mathbb{Q})$ unter Verwendung von $r$ Klassen gleichzeitig beschränkt und etabliert:
\begin{equation}
\rank E(\mathbb{Q}) \leq \ord_{s=1} L(E,s)\,.
\end{equation}

\textit{Status:} Für $r = 1$ ist dies Kolyvagins Methode. Für allgemeines $r$ haben Mazur-Rubin \cite{MazurRubin2004} die Theorie der Kolyvagin-Systeme abstrakt entwickelt, und Burns-Sano \cite{BurnsSano2021} sowie Burns-Sakamoto-Sano \cite{BurnsSakamotoSano2025} haben diese auf höherrangige Euler-, Kolyvagin- und Stark-Systeme über allgemeinen Koeffizientenringen erweitert. Kim \cite{Kim2024} hat eine Kolyvagin-System-theoretische Verfeinerung der Gross-Zagier-Formel erhalten. Die \textit{Existenz} eines Kolyvagin-Systems hinreichender Tiefe für beliebigen Rang ist jedoch nicht etabliert.

\textbf{(H3) $\Sha$-Endlichkeit für allgemeinen Rang.} Axiom~I muss unbedingt gelten:
\begin{equation}
|\Sha(E/\mathbb{Q})| < \infty \quad \forall\, E/\mathbb{Q}\,.
\end{equation}

\textit{Status:} Offen für $\rank \geq 2$. Dies wird allgemein erwartet, aber es existiert kein allgemeiner Beweis.

\subsection{Die Positivitätslücke}

Die Situation kann zusammengefasst werden:

\begin{center}
\begin{tabular}{lccc}
\toprule
Axiom & Rang 0 & Rang 1 & Rang $\geq 2$ \\
\midrule
I (Endl.\ $\Sha$) & $\checkmark$ & $\checkmark$ & Offen \\
II (Höhen-Pos.) & $\checkmark$ & $\checkmark$ & $\checkmark$ \\
III (Iwasawa) & $\checkmark$ & $\checkmark$ & Teilweise \\
IV (Euler-Sys.) & $\checkmark$ & $\checkmark$ & Offen \\
\bottomrule
\end{tabular}
\end{center}

Die ,,Positivität'' (Axiom~II) ist immer verfügbar -- sie ist ein Theorem. Die Obstruktion liegt in der ,,Komposition'' (Axiome~I, III, IV): die Euler-System-/Kolyvagin-/Iwasawa-Maschinerie ist für höheren Rang noch nicht vollständig entwickelt.

Zusammenfassend: die Positivität (Axiom~II) ist etabliert, aber das Kompositionsgesetz (Axiome~I, III, IV) ist für Rang~$\geq 2$ unvollständig. Die BSD-Formel wird erst dann zur erwarteten Normalform, wenn die Euler-System-Komposition, die Höhen-/Zykelkopplung und die Sha-Kontrolle für den betreffenden Rang etabliert sind. Diese Arbeit ist Teil eines breiteren Programms \cite{Geiger2026e}; die Verbindungen sind rein motivational und werden in keinem Beweis verwendet.


% ===================================================================
% 7. IWASAWA-THEORIE UND KOMPOSITION
% ===================================================================
\section{Iwasawa-Theorie als Kompositionskohärenz}
\label{sec:iwasawa}

\subsection{\texorpdfstring{Die $p$-adische Perspektive}{Die p-adische Perspektive}}

Die Iwasawa-Theorie liefert die klarste Realisierung von Axiom~III (monotone Kontrolle) und verbindet es natürlich mit Axiom~IV (Euler-System-Kohärenz). Die Hauptvermutung für elliptische Kurven \cite{Kato2004,SkinnerUrban2014} ist in den folgenden spezifischen Fällen etabliert:
\begin{itemize}
\item \textbf{Katos Teilbarkeit (eine Richtung):} Für alle Primzahlen $p$, bei denen $E$ gute ordinäre Reduktion hat und $\bar{\rho}_{E,p}$ surjektiv ist. Gilt für \textit{alle Ränge}, liefert aber nur die Inklusion $\mathrm{char}_\Lambda(\Sel^\vee) \subseteq (\mathcal{L}_p)$, d.h.\ das charakteristische Ideal des Selmer-Duals ist im vom $p$-adischen $L$-Funktion erzeugten Ideal enthalten, was eine obere Schranke für den Selmer-Modul liefert.
\item \textbf{Skinner-Urban (volle Gleichheit):} Für ordinäre Primzahlen $p$ mit $p \nmid a_p$, unter der Voraussetzung der Irreduzibilität von $\bar{\rho}_{E,p}$ und weiterer lokaler Hypothesen. Deckt \textit{keine} supersingularen Primzahlen ab.
\item \textbf{Supersingulare Primzahlen:} Erfordern die vorzeichenbehaftete Selmer-Gruppen-Formulierung (Kobayashi, Pollack, Sprung); die Hauptvermutung in diesem Rahmen bleibt weitgehend offen.
\item \textbf{Nicht abgedeckt:} Primzahlen schlechter additiver Reduktion, Primzahlen bei denen $\bar{\rho}_{E,p}$ reduzibel ist, und der allgemeine Fall ohne residuale Irreduzibilität.
\end{itemize}
Die Hauptvermutung behauptet:
\begin{equation}
\mathrm{char}_{\Lambda}\!\left(\Sel_{p^\infty}(E/\mathbb{Q}_\infty)^\vee\right) = (L_p(E))\,,
\label{eq:iwasawa_mc}
\end{equation}
wobei $\Lambda = \mathbb{Z}_p[[T]]$ die Iwasawa-Algebra, $\Sel_{p^\infty}(E/\mathbb{Q}_\infty)^\vee$ der Pontryagin-Dual der Selmer-Gruppe über der zyklotomischen $\mathbb{Z}_p$-Erweiterung und $L_p(E)$ die $p$-adische $L$-Funktion ist.

\begin{proposition}[Iwasawa-Hauptvermutung als Kompositionsgesetz]
\label{prop:iwasawa_composition}
Die Iwasawa-Hauptvermutung liefert die Kompositionsstruktur im folgenden präzisen Sinne: sei $\mathrm{char}_\Lambda(M)$ das charakteristische Ideal eines endlich erzeugten Torsions-$\Lambda$-Moduls $M$. Die Hauptvermutung identifiziert $\mathrm{char}_\Lambda(\Sel_{p^\infty}(E/\mathbb{Q}_\infty)^\vee) = (L_p(E))$. Die ,,Komposition'' ist das Kontrolltheorem (Mazur): bei Spezialisierung auf Schicht $n$ des $\mathbb{Z}_p$-Turms ist die Selmer-Gruppe auf Stufe $n$ mit der Selmer-Gruppe auf Stufe $n-1$ durch eine exakte Sequenz verbunden, deren Kern und Kokern durch den Euler-Faktor bei $p$ kontrolliert werden. Diese schichtweise Kompatibilität ist das arithmetische Analogon der Norm-Kompatibilität in Euler-Systemen.
\end{proposition}

Die algebraische Struktur ist: $\Lambda = \mathbb{Z}_p[[T]]$ ist ein Potenzreihenring, und das charakteristische Ideal wird von einem einzigen Element $f_E(T)$ erzeugt. Die Spezialisierungsabbildungen $\chi_n: \Lambda \to \mathbb{Z}_p[\Gal(\mathbb{Q}_n/\mathbb{Q})]$, die durch Auswertung bei $T = \zeta_{p^n} - 1$ erhalten werden, bilden ein projektives System:
\begin{equation}
\chi_m = \chi_n \circ \mathrm{pr}_{m,n} \quad \text{für } m \geq n\,,
\end{equation}
wobei $\mathrm{pr}_{m,n}$ die natürliche Projektion ist. Das Kontrolltheorem stellt sicher, dass $\Sel_{p^\infty}(E/\mathbb{Q}_n)$ durch $\Sel_{p^\infty}(E/\mathbb{Q}_\infty)$ via Restriktion mit beschränktem Kern und Kokern bestimmt wird. Dieses kontrollierte kompatible System ist das $p$-adische Kompositionsgesetz, das die Arithmetik von $E$ über den Turm hinweg einschränkt.

\subsection{\texorpdfstring{Die $\mu = 0$-Vermutung}{Die mu = 0-Vermutung}}

Die Iwasawa-$\mu$-Invariante misst die ,,$p$-adische Masse'' von $\Sha$ im Turm. Die Vermutung $\mu = 0$ (bewiesen von Ferrero-Washington für abelsche Körper, allgemein erwartet) stellt sicher, dass kein unendlicher $p$-Potenzbeitrag in $\Sha$ verborgen ist. Dies ist präzise Axiom~I eingeschränkt auf den $p$-adischen Turm.

\subsection{\texorpdfstring{Spezialisierung auf $s = 1$}{Spezialisierung auf s = 1}}

Die Schlüsselverbindung zwischen Iwasawa-Theorie und BSD ist die \textit{Spezialisierung}: die Auswertung der $p$-adischen $L$-Funktion bei $T = 0$ (oder ihrer Ableitungen) gewinnt die komplexen $L$-Funktionswerte zurück:
\begin{equation}
\begin{aligned}
L_p(E, 0) &\sim L(E, 1)/\Omega_E^+ \\
&\quad (\text{bis auf Perioden und Euler-Faktoren})\,.
\end{aligned}
\end{equation}

Für Rang $r$ betrifft die relevante Spezialisierung die $r$-te Ableitung von $L_p(E, T)$ bei $T = 0$. Die Iwasawa-Hauptvermutung bezieht diese dann auf das charakteristische Ideal des Selmer-Moduls und liefert eine $p$-adische Version der BSD-Formel.

\subsection{\texorpdfstring{Die $p$-adische zu archimedische Brücke}{Die p-adische zu archimedische Brücke}}
\label{sec:p-adic-bridge}

Die Iwasawa-Hauptvermutung und die Euler-System-Maschinerie (Axiome~III--IV) leben natürlich in der $p$-adischen Welt, während die BSD-Vermutung die \textit{komplexe} $L$-Funktion $L(E,s)$ an der archimedischen Stelle $s = 1$ betrifft. Der Übergang zwischen diesen beiden Domänen wird durch Fontaines Vergleichsisomorphismus $H^1_{\mathrm{dR}}(E/\mathbb{Q}_p) \cong H^1_{\mathrm{\acute{e}t}}(E, \mathbb{Q}_p) \otimes B_{\mathrm{dR}}$ aus der $p$-adischen Hodge-Theorie kontrolliert, der ein Wörterbuch liefert zwischen $p$-adischen Perioden und komplexen Perioden ($\Omega_E$), zwischen $p$-adischen Höhen ($\hat{h}_p$) und archimedischen Höhen ($\hat{h}$), sowie zwischen $p$-adischen $L$-Werten ($L_p(E,0)$) und komplexen $L$-Werten ($L(E,1)$).

Für Rang~0 ist die Brücke gut verstanden: bei einer Primzahl $p$ guter ordinärer Reduktion ergibt die Interpolationsformel $L_p(E,0) = (1 - \alpha_p^{-1})^2 \cdot L(E,1)/\Omega_E^+$ bis auf eine $p$-adische Einheit, wobei $\alpha_p$ die Einheitswurzel von $X^2 - a_p X + p$ ist (für multiplikative Reduktion führt die Mazur-Tate-Teitelbaum-Vermutung, bewiesen von Greenberg-Stevens \cite{GreenbergStevens1993}, die $\mathcal{L}$-Invariante ein). Für Rang~1 liefern die Gross-Zagier- und Perrin-Riou-Formeln \cite{PerrinRiou1987} kompatible archimedische und $p$-adische Höhenkopplungen und schließen die Brücke auf der Ebene der ersten Ableitungen.

Für Rang $\geq 2$ ist die Brücke die tiefste Obstruktion im Positivitäts-Framework. Selbst wenn eine $p$-adische Höhenkopplung $L_p^{(r)}(E,0) \sim \det(\hat{h}_p(P_i, P_j))$ über diagonale Zykel \cite{DarmonRotger2017} etabliert würde, erfordert die Übersetzung zur archimedischen Formel $L^{(r)}(E,1) \sim R_E$ zwei zusätzliche Zutaten: (i) einen Vergleich zwischen $p$-adischen und archimedischen Regulatoren, $R_{E,p} \sim R_E$, der bis auf eine $p$-adische Einheit für CM-Kurven bekannt, aber im Allgemeinen offen ist; und (ii) eine Interpolationsformel für die $r$-te Ableitung, die Mazur-Tate-Teitelbaum über Rang~0 hinaus verallgemeinert.

Diese Brücke kristallisiert die strukturelle Spannung in unserem Framework: die \textit{Positivität} (Axiom~II) ist ein archimedisches Phänomen -- $\hat{h}$ ist eine reellwertige positiv-definite Form -- während die \textit{Komposition} (Axiome~III--IV) ein $p$-adisches Phänomen ist -- Euler-Systeme und Iwasawa-Moduln leben über $\mathbb{Z}_p$. Die Vereinigung dieser beiden ,,Positivitätsdomänen'' über den Vergleichsisomorphismus ist die zentrale technische Herausforderung für die Erweiterung des BSD-Positivitätsprogramms jenseits von Rang~1.


% ===================================================================
% 8. DIE PRAEZISE OBSTRUKTION
% ===================================================================
\section{Die präzise Obstruktion}
\label{sec:obstruction}

\subsection{Was einen vollständigen Beweis blockiert}

Das Positivitäts-Framework identifiziert einen präzisen Engpass: das \textit{Fehlen einer universellen höheren Gross-Zagier-Formel}. Genauer:

\begin{enumerate}
\item \textbf{Für Rang 1:} Der Heegner-Punkt $P_K$ liefert eine kanonische ,,geometrische Richtung'', deren Höhe an $L'(E,1)$ koppelt. Das Kompositionsgesetz (Kolyvagins Euler-System) schließt dann Alternativen aus.

\item \textbf{Für Rang $r$:} Man benötigt $r$ unabhängige kanonische Punkte $P_1, \ldots, P_r$, deren Höhen\textit{matrix} an $L^{(r)}(E,1)$ koppelt. Die Kandidaten sind:
\begin{itemize}
\item Höhere Heegner-Punkte auf Shimura-Varietäten der Dimension $> 1$.
\item Verallgemeinerte Heegner-Zykel (Bertolini-Darmon-Prasanna \cite{BDP2013}).
\item Diagonalrestriktion $p$-adischer Familien (Darmon-Rotger \cite{DarmonRotger2017}).
\end{itemize}
Keine liefert derzeit eine vollständige ,,höhere Gross-Zagier-Formel'' der Form~\eqref{eq:higher_gz}.
\end{enumerate}

\subsection{Die Obstruktion in Positivitätssprache}

Die Obstruktion kann präzise formuliert werden:

\begin{quote}
\textit{Die Positivität von $R_E = \det(\langle P_i, P_j \rangle) > 0$ ist ein Theorem (aus der Positiv-Definitheit von $\hat{h}$). Das fehlende Stück ist eine Formel, die $L^{(r)}(E,1)$ gleich $R_E$ mal berechenbare Konstanten setzt -- d.h., die explizite Kopplung der analytischen und geometrischen positiven Strukturen.}
\end{quote}

Sobald eine solche Kopplung etabliert ist, schließt sich das No-Go-Argument:
\begin{enumerate}
\item $L^{(r)}(E,1) = C_r \cdot R_E > 0$ (Positivität der Höhen), also $\ord_{s=1} L(E,s) \leq r$.
\item Die Euler-System-Schranke (Axiom~IV) ergibt $\rank E(\mathbb{Q}) \leq \ord_{s=1} L(E,s)$, also $\ord_{s=1} L(E,s) \geq r$.
\item Zusammen: $\ord_{s=1} L(E,s) = r = \rank E(\mathbb{Q})$.
\end{enumerate}

Das gesamte Argument beruht auf der Kopplung zweier positiver Strukturen -- einer analytischen ($L$-Ableitungen), einer geometrischen (Höhenpaarung) -- und dem Nachweis, dass sie bei derselben Ordnung ,,sättigen'' müssen.


% ===================================================================
% 9. VERBINDUNG ZUM BREITEREN PROGRAMM
% ===================================================================
\section{Verbindung zum breiteren Programm}
\label{sec:connection}

\subsection{Das Positivitätsmuster über die Mathematik}

Das BSD-Framework fügt sich in dasselbe Positivitätsmuster ein, das beobachtet wird in:

\begin{center}
\footnotesize
\resizebox{\columnwidth}{!}{%
\begin{tabular}{@{}llll@{}}
\toprule
Problem & Positive Form & Komposition & Status \\
\midrule
Weil-Verm. & Frobenius-EW & \'Etale Koh. & Bewiesen \\
Hodge-Verm. & HR-Bilinear & Galois-HR & Offen \\
RH \cite{GeigerRH2026c} & Spektr.\ Op. & Resolvente & Offen \\
CRM/QG \cite{FST-Unified2026} & Kapazität & RG-Fluss & Sättigungs-Thm \\
\textbf{BSD} & \textbf{N\'eron-Tate} & \textbf{Euler-Sys.} & \textbf{Rang $\leq 1$} \\
\bottomrule
\end{tabular}
}
\end{center}
\noindent{\footnotesize CRM/QG = Curvature Relaxation Model; \href{https://doi.org/10.5281/zenodo.18728935}{doi:10.5281/zenodo.18728935}.}

In jedem Fall:
\begin{itemize}
\item Die \textit{Positivität} ist der ,,leichte'' Teil (ein Theorem oder eine gut motivierte Bedingung).
\item Die \textit{Komposition/Kohärenz} ist der ,,schwere'' Teil (erfordert tiefe strukturelle Arbeit).
\item Die \textit{Konjunktion} von Positivität und Komposition erzwingt das gewünschte Resultat.
\end{itemize}

\subsection{Verbindung zum Hodge-Ansatz}

Die N\'eron-Tate-Höhe $\hat{h}$ auf $E(\mathbb{Q})$ und die Hodge-Riemann-Form $Q_L$ auf $H^{p,p}(X)$ spielen strukturell identische Rollen:
\begin{itemize}
\item Beide sind kanonische positiv-definite Bilinearformen.
\item Beide detektieren ,,arithmetische Richtungen'' (rationale Punkte / algebraische Zykel).
\item Beide erfordern ein ,,Kompositionsgesetz'' (Euler-Systeme / Galois-HR-Stabilität), bevor die gewünschte Schlussfolgerung als Theorem erreichbar wird.
\end{itemize}

Das Hodge-Analogon der Gross-Zagier-Formel wäre: eine explizite Identität, die die ,,analytische'' Hodge-Riemann-Form an die ,,geometrische'' Zykelklassenabbildung koppelt und zeigt, dass Positivität der einen Algebraizität der anderen erzwingt.

\subsection{Verbindung zur Szpiro-Vermutung}
\label{sec:szpiro}

Das BSD-Framework verbindet sich über die Tamagawa-Zahlen $c_p$ in der Leitterm-Formel~\eqref{eq:bsd_formula} mit der Szpiro-Vermutung (äquivalent zur $abc$-Vermutung).

\subsubsection{Szpiro als lokale Diskriminantenexponenten-Schranke}

\begin{conjecture}[Szpiro {\cite{Szpiro1981}}]
\label{conj:szpiro}
Für jedes $\varepsilon > 0$ existiert $C_\varepsilon > 0$, sodass für jede elliptische Kurve $E/\mathbb{Q}$ mit minimaler Diskriminante $\Delta_{\min}$ und Führer $N$ gilt:
\begin{equation}
\log|\Delta_{\min}| \leq (6+\varepsilon)\log N + C_\varepsilon\,.
\label{eq:szpiro}
\end{equation}
\end{conjecture}

Für eine semistabile Kurve $E/\mathbb{Q}$ ist jede schlechte Reduktion multiplikativ. Schreibe
\[
  n_p := v_p(\Delta_{\min}) .
\]
Dann ist der Führer $N=\prod_{p\mid\Delta}p$ und
\begin{equation}
\sum_{p\mid N} n_p\log p
= \log|\Delta_{\min}|
\leq (6+\varepsilon)\sum_{p\mid N}\log p + C_\varepsilon\,.
\label{eq:szpiro_tamagawa}
\end{equation}
Bei split-multiplikativer Reduktion an $p$ ist die Tamagawa-Zahl gleich diesem lokalen Exponenten, $c_p=n_p$ (\cite{Silverman2009}, Tabelle~IV.9.1). Bei nonsplit-multiplikativer Reduktion kann $c_p$ dagegen nur $1$ oder $2$ sein, obwohl $n_p$ größer ist. Szpiro ist daher grundsätzlich eine \emph{primgewichtete lokale Diskriminantenexponenten-Schranke}; nur im split-multiplikativen Spezialfall kann sie wortwörtlich als primgewichtete Tamagawa-Exponenten-Schranke gelesen werden. Dies ist stärker und strukturierter als eine ungewichtete Schranke für $\prod_p c_p$ (oder sogar für $\prod_p n_p$): Ein grosser Exponent an einer grossen Primzahl trägt mehr Diskriminantenmasse als derselbe Exponent an einer kleinen Primzahl.

\subsubsection{Diagnosekette: BSD plus Schranken zeigt den fehlenden gewichteten Schritt}

Die BSD-Formel~\eqref{eq:bsd_formula} liefert:
\begin{equation}
\prod_p c_p = \frac{L^{(r)}(E,1)}{r!}\cdot
\frac{|E_{\mathrm{tors}}|^2}{\Omega_E\cdot R_E\cdot |\Sha|}\,.
\label{eq:tamagawa_from_bsd}
\end{equation}
Diese Identität kontrolliert das \emph{ungewichtete Produkt} der Tamagawa-Faktoren. Sie impliziert für sich genommen nicht Szpiro, weil Szpiro die primgewichtete Summe in~\eqref{eq:szpiro_tamagawa} kontrolliert. Eine Route von BSD zu Szpiro braucht daher neben den folgenden Zutaten ein zusätzliches gewichtetes lokales Upgrade:
\begin{enumerate}
\item[\textbf{(B1)}] \textbf{Obere Schranke für $L$-Werte:} $L^{(r)}(E,1)/r! \leq C\cdot N^{r/2+\varepsilon}$. Unbedingt für $r=0$ (Konvexitätsschranke); für allgemeines $r$ bedingt auf GRH.
\item[\textbf{(B2)}] \textbf{Untere Periodenschranke:} $\Omega_E \geq c\cdot N^{-1/2-\varepsilon}$. Dies folgt aus Schranken für den modularen Parametrisierungsgrad.
\item[\textbf{(B3)}] \textbf{Untere Regulator-Schranke:} $R_E \geq c\cdot(\log N)^r$. Dies verlangt eine effektive untere Schranke für die N\'eron--Tate-Höhe nicht-torsioneller Punkte.
\item[\textbf{(B4)}] \textbf{Torsionsschranke:} $|E_{\mathrm{tors}}|^2\leq 256$. Dies ist Mazurs Torsionstheorem ($|E(\mathbb{Q})_{\mathrm{tors}}|\leq 16$).
\item[\textbf{(B5)}] \textbf{Gewichtetes lokales Upgrade:} ein Mechanismus, der ungewichtete Kontrolle von $\prod_p c_p$ in Kontrolle der tatsächlichen lokalen Exponenten $\sum_{p\mid N}n_p\log p$ übersetzt, einschließlich der split-/nonsplit-multiplikativen Lücke, oder direkt jeden lokalen Diskriminantenbeitrag durch Führerdaten beschränkt.
\end{enumerate}
Aus (B1)--(B4) unter BSD folgt nur:
\begin{equation}
\prod_p c_p \leq C'\cdot\frac{N^{r/2+1/2+2\varepsilon}}{(\log N)^r}\,.
\label{eq:tamagawa_bound}
\end{equation}
Für semistabile Kurven vom Rang~$0$ ist dies eine starke ungewichtete Tamagawa-Produkt-Beschränkung. Eine Szpiro-Schranke wird daraus erst, wenn (B5) sowohl die fehlenden Primgewichte als auch den lokalen Übergang von Tamagawa-Faktoren zu den tatsächlichen Diskriminantenexponenten $n_p$ liefert.

\begin{remark}[Zirkularität in Schritt (B3)]
\label{rem:szpiro_circularity}
Schon vor dem gewichteten Upgrade hat die Kette eine zweite Obstruktion: Zirkularität in (B3). Die beste bekannte effektive untere Höhenschranke ist Silvermans $\hat{h}(P)\geq c\cdot\log N$ für nicht-torsionelle Punkte $P$~\cite{Silverman2009}; sie ist \textit{bedingt auf die $abc$-Vermutung}. Da $abc\Leftrightarrow$ Szpiro, wäre die logische Struktur:
\[
\begin{gathered}
\text{BSD}+\underbrace{\text{abc}}_{\text{in (B3) genutzt}}\\
+\underbrace{\text{gewichtetes lokales Upgrade}}_{\text{fehlendes (B5)}}
\Longrightarrow
\underbrace{\text{Szpiro}}_{\Leftrightarrow\,\text{abc}}\,.
\end{gathered}
\]
Um diese Zirkularität zu brechen, braucht man eine Höhenschranke, die nicht von der Diskriminante oder von $abc$ abhängt, und einen separaten Mechanismus für die primgewichteten lokalen Exponenten.
\end{remark}

\subsubsection{Eine Theta-Produkt-Brücke}

Die folgende Identität verbindet Theta-Werte an Torsionspunkten mit der Dedekind-Eta-Funktion und damit mit der modularen Diskriminante. Sie liefert eine mögliche nicht-zirkuläre Verbindung zwischen Torsionsstruktur und Diskriminante.

% TOOL-NOTE: Die Theta-Eta-Produktidentität ist konsolidiert in
%   FST_MATHEMATICS/_tools/THETA_ETA_IDENTITY.md
% Identische Aussage und Notation. Dieselbe Identität wird in
% abc/abc_Theta_Tamagawa_EN.tex Proposition 4.1 verwendet.

\begin{proposition}[Theta-Eta-Produktidentität an $l$-Torsion]
\label{prop:theta_eta}
Für jede ungerade Primzahl $l$ und $\tau\in\mathcal{H}$, mit $q=e^{2\pi i\tau}$, gilt:
\begin{equation}
\prod_{j=1}^{l-1}\theta_1\!\left(\frac{j}{l}\,\bigg|\,\tau\right)
= l\cdot\eta(\tau)^{l-3}\cdot\eta(l\tau)^2\,,
\label{eq:theta_eta_identity}
\end{equation}
wobei $\theta_1(z|\tau)=-i\sum_{n\in\mathbb{Z}}(-1)^nq^{(n+1/2)^2/2}e^{(2n+1)\pi iz}$ die Jacobi-Theta-Funktion und $\eta(\tau)=q^{1/24}\prod_{n=1}^{\infty}(1-q^n)$ die Dedekind-Eta-Funktion ist.
\end{proposition}

\begin{proof}
Das Jacobi-Tripelprodukt liefert $\theta_1(z|\tau)=-2q^{1/8}\sin(\pi z)\prod_{n=1}^{\infty}(1-q^n)(1-e^{2\pi iz}q^n)(1-e^{-2\pi iz}q^n)$.
Setzt man $\omega=e^{2\pi i/l}$, dann ergeben die zyklotomische Identität $\prod_{j=1}^{l-1}(1-\omega^jx)=1+x+\cdots+x^{l-1}$ für jeden Faktor des unendlichen Produkts und die Sinus-Identität $\prod_{j=1}^{l-1}\sin(\pi j/l)=l/2^{l-1}$:
\[
\prod_{j=1}^{l-1}\theta_1(j/l|\tau)
= l\,q^{(l-1)/8}
\prod_{n=1}^{\infty}(1-q^n)^{l-3}(1-q^{ln})^2\,.
\]
Da $\eta(\tau)=q^{1/24}\prod_{n\geq1}(1-q^n)$ und $\eta(l\tau)=q^{l/24}\prod_{n\geq1}(1-q^{ln})$ gilt, besitzt das Eta-Produkt $l\eta(\tau)^{l-3}\eta(l\tau)^2$ genau den Vorfaktor $lq^{(l-3)/24+2l/24}=lq^{(l-1)/8}$ und dasselbe unendliche Produkt. Damit ist die Identität bewiesen.
\end{proof}

\begin{remark}[Verbindung zur Diskriminante]
\label{rem:theta_discriminant}
Da $\eta(\tau)^{24}=\Delta(\tau)/(2\pi)^{12}$ gilt, wobei $\Delta$ die modulare Diskriminante ist, folgt aus~\eqref{eq:theta_eta_identity}:
\begin{equation}
\begin{aligned}
\sum_{j=1}^{l-1}\log|\theta_1(j/l|\tau)|
&= \log l
 + \frac{l-3}{24}\log\frac{|\Delta(\tau)|}{(2\pi)^{12}} \\
&\quad + \frac{2}{24}\log\frac{|\Delta(l\tau)|}{(2\pi)^{12}}\,.
\end{aligned}
\label{eq:theta_log_discriminant}
\end{equation}
Für grosses $l$ dominiert der Term $(l/24)\log|\Delta(\tau)|$: Das Produkt der Theta-Werte an $l$-Torsionspunkten kodiert die Diskriminante.
\end{remark}

\begin{remark}[Tate-Kurve und Szpiro-Verbindung]
\label{rem:tate_szpiro}
Für eine elliptische Kurve $E/\mathbb{Q}_p$ mit split-multiplikativer Reduktion (Tate-Kurve $E_q\cong\mathbb{G}_m/q^{\mathbb{Z}}$) entsprechen die $l$-Torsionspunkte den Klassen $\zeta_l^j\cdot q^{k/l}$ für $0\leq j,k<l$. Die Identität~\eqref{eq:theta_eta_identity} verbindet die $l$-Torsionsstruktur von $E$ mit dem lokalen Diskriminantenexponenten $n_p=v_p(\Delta_{\min})$ (äquivalent zu $c_p$ an split-multiplikativen Stellen), \textit{ohne Höhenschranken oder die $abc$-Vermutung zu verwenden}. Dies weist auf einen möglichen nicht-zirkulären Pfad zu Szpiro hin: Wenn man die $l$-Torsions-Theta-Werte in Termen des Führers $N$ beschränken kann, etwa über die Galois-Darstellung $\bar{\rho}_{E,l}$ oder die explizite Weil-Formel für $L(E,s)$, dann würde die daraus folgende Schranke für $n_p$ direkt die lokale Eingabe für die Szpiro-Vermutung liefern.

Der fehlende Schritt ist eine primgewichtete lokale Führer-Schranke der Form:
\begin{equation}
\begin{aligned}
\sum_{p\mid N}\sum_{j=1}^{l-1}
-\log|\theta_1(j/l|\tau_p)|_p
&\leq C(l,\varepsilon)\cdot\log N\,,
\end{aligned}
\label{eq:theta_conductor_bound}
\end{equation}
die derzeit nicht etabliert ist, aber das richtige gewichtete Ziel formuliert. Die Beschränkung der Theta-Werte über den Führer verbindet sich natürlich mit der Galois-Wirkung auf $E[l]$, wobei Serres Modularitätsresultate und Level-Lowering die Verbindung zwischen $l$-Torsionsdarstellung und Führer liefern.
\end{remark}

\begin{remark}[Bezug zu Mochizukis IUT-Ansatz]
\label{rem:iut_comparison}
Die Theta-Eta-Produktidentität~\eqref{eq:theta_eta_identity} operiert auf denselben Objekten -- Theta-Werten an Torsionspunkten von Tate-Kurven -- die auch in Mochizukis Inter-Universal Teichmuller Theory (IUT) auftreten. Der IUT-Ansatz verwendet jedoch Theta-Werte $q^{j^2}$, also quadratische Abhängigkeit von $j$, die aus der multiradialen Darstellung stammt, während die Identität~\eqref{eq:theta_eta_identity} $\theta_1(j/l|\tau)$ mit linearer Abhängigkeit von $j$ betrifft. Diese Unterscheidung wird in~\cite{GeigerIUT2026} analysiert; dort wird gezeigt, dass der $O(l)$-gegen-$O(l^2)$-Skalierungsunterschied zwischen linearen und quadratischen Theta-Werten eine strukturelle Obstruktion gegen die Reproduktion der IUT-Schranke mit klassischen Methoden bildet. Der vorliegende Ansatz unterscheidet sich grundlegend: Er versucht nicht, den IUT-Mechanismus zu reproduzieren, sondern nutzt die Produktidentität als direkte Brücke von Torsion zu Diskriminante innerhalb des BSD-Frameworks.
\end{remark}


% ===================================================================
% 10. DISKUSSION
% ===================================================================
\section{Diskussion}
\label{sec:discussion}

\subsection{Was wurde erreicht}

\begin{enumerate}
\item \textbf{Axiomatisierung.} Vier Axiome -- endliche Kapazität, Höhenpositivität, kontrolliertes Iwasawa-Wachstum, Euler-System-Komposition -- die die BSD-Formel bedingt als erwartete Produktnormalform identifizieren, sobald (H1)--(H3) geliefert sind.

\item \textbf{Rang $\leq 1$ als abgeschlossene Instanz.} Das Gross-Zagier + Kolyvagin-Argument ist natürlich ein Positivitäts-No-Go-Beweis: es koppelt eine analytische Ableitung an eine nichtnegative Höhe und verwendet Komposition (Euler-Systeme) zum Ausschluss von Alternativen.

\item \textbf{Höherer Rang identifiziert.} Die Obstruktion für höheren Rang ist präzise das Fehlen von (H1) einer universellen höheren Gross-Zagier-Formel und (H2) eines vollständigen Kolyvagin-Systems -- d.h., das Kompositionsgesetz ist unvollständig. Die Kolyvagin--Euler-System-Struktur wurde in analoger Form auf das Begleitpaper zu P vs NP übertragen, wo eine Hierarchie von Zeugen rechnerische Härte nach oben propagiert -- in Spiegelung der arithmetischen Propagation von Nichttrivialität im BSD-Kontext.

\item \textbf{Ausschluss von Alternativen.} Kein alternatives Führungstermprofil ist mit allen vier Axiomen kompatibel, \textit{sofern} eine Kopplung zwischen $L^{(r)}(E,1)$ und der Höhenpaarung existiert.
\end{enumerate}

\subsection{Einschränkungen}

\begin{enumerate}
\item \textbf{Die Arbeit beweist BSD nicht.} Sie formuliert es als Positivitäts-Normalform-Aussage um und identifiziert die fehlenden Brückenmodule.

\item \textbf{Axiom~I ($\Sha$-Endlichkeit) ist für Rang $\geq 2$ konjektural.} Dies ist selbst ein großes offenes Problem, obwohl allgemein erwartet.

\item \textbf{Das ,,höhere Gross-Zagier'' ist spekulativ.} Keine vollständige Formel der Form~\eqref{eq:higher_gz} existiert für $r \geq 2$. Jüngerer Fortschritt (Darmon-Rotger, Loeffler-Skinner-Zerbes) ist vielversprechend, aber unvollständig.

\item \textbf{Das axiomatische Framework ist deskriptiv, nicht konstruktiv.} Die Axiome identifizieren, was wahr sein muss, nicht wie man es beweist.
\end{enumerate}

\subsection{Ausblick}

\begin{enumerate}
\item \textbf{$p$-adische BSD über Iwasawa-Theorie.} Der stärkste kurzfristige Fortschritt kommt wahrscheinlich aus der Erweiterung der Iwasawa-Hauptvermutung auf höherrangige Situationen unter Verwendung der Euler-Systeme von Loeffler-Skinner-Zerbes \cite{LSZ2020} und der höherrangigen Kolyvagin-Stark-System-Theorie von Burns-Sakamoto-Sano \cite{BurnsSano2021,BurnsSakamotoSano2025}.

\item \textbf{Höheres Gross-Zagier.} Darmons Programm der ,,Stark-Heegner-Punkte'' und Diagonalzykel liefert Kandidatenformeln für Rang 2. Deren rigorose Etablierung würde die Positivitätsschleife für $r = 2$ schließen.

\item \textbf{Numerische Verifikation.} Die BSD-Formel wurde numerisch für Tausende von Kurven des Rangs $\leq 3$ verifiziert (siehe \cite{LMFDB}). Die Erweiterung dieser Berechnungen auf höhere Präzision und höheren Rang testet die Vorhersagen des Frameworks.

\item \textbf{Statistische BSD.} Die Arbeit von Bhargava-Shankar \cite{BhargavaShankar2015} zeigt, dass ein positiver Anteil elliptischer Kurven über $\mathbb{Q}$ Rang~0 hat und die BSD-Vermutung erfüllt. Diese statistischen Resultate liefern weitere Evidenz dafür, dass das Positivitäts-Framework das generische Verhalten erfasst.
\end{enumerate}


\begin{remark}[Arakelov-Brücke zu Hodge]
\label{rem:arakelov-bridge}
Die N\'eron--Tate-Höhenpaarung kodiert Positivität an endlichen (nichtarchimedischen) Primstellen. Zusammen mit der Hodge--Riemann-Form (Positivität an der unendlichen Primstelle, vgl.\ das Begleitpaper zu Hodge) liefert die Arakelov-Schnittpaarung den natürlichen vereinheitlichten Rahmen, in dem beide Spezialfälle eines einzigen Positivitätsprinzips sind.
\end{remark}


\begin{remark}[Geometrische Interpretation des Regulators]
\label{rem:entropic-capacity}
Der Regulator $R_E = \det(\langle P_i, P_j \rangle_{\mathrm{NT}})$ ist das Kovolumen des Mordell--Weil-Gitters $E(\mathbb{Q})/E(\mathbb{Q})_{\mathrm{tors}}$ in der N\'eron--Tate-Metrik. Die BSD-Formel setzt das analytische Datum $L^{(r)}(E,1)$ mit diesem geometrischen Volumen gleich (mal lokale Korrekturfaktoren). Eine Verbindung zur Zufallsmatrizentheorie besteht über die Keating--Snaith-Vermutungen zur Verteilung von $L$-Funktions-Ableitungen, aber wir entwickeln dies hier nicht weiter.
\end{remark}

\subsection{Forschungsrichtungen für höheren Rang}
\label{sec:higher_rank_directions}

Die folgenden Bemerkungen untersuchen konkrete Ansätze zur Herstellung der fehlenden Hypothesen (H1)--(H3) für Rang~$\geq 2$.

\subsubsection{Diagonale Zykel und höhere Gross-Zagier-Formeln}

\begin{remark}[Diagonal-Zyklus / Euler-System-Hybrid für Rang-2 CM-Kurven]
\label{rem:diagonal-euler}
Für eine CM-Kurve $E/\mathbb{Q}$ mit Rang~2 könnte man zwei unabhängige Kohomologieklassen -- eine von einem Heegner-Zyklus und eine von einem Beilinson-Flach-Element -- konstruieren, deren Höhen die volle Regulatormatrix aufspannen. Darmon und Rotger~\cite{DarmonRotger2017} liefern Teilresultate über diagonale Zykel auf dreifachen Produkten modularer Kurven.
\end{remark}

\begin{remark}[Heegner-Simplex-Ansatz für Rang $r$]
\label{rem:heegner-simplex}
Für Rang~$r$ sucht man $r$-dimensionale Zykel auf Kuga-Sato-Varietäten oder höherdimensionalen Shimura-Varietäten, deren Abel-Jacobi-Bilder eine Untergruppe vollen Ranges der relevanten Selmer-Gruppe erzeugen. Das gruppentheoretische Framework (Ersetzung von $\mathrm{GL}_2$ durch $\mathrm{GSp}_{2r}$) liefern Loeffler--Skinner--Zerbes~\cite{LSZ2020}, aber die explizite Höhenformel bleibt offen.
\end{remark}

\begin{remark}[Plektische Gross-Zagier-Formel]
\label{rem:plectic-gz}
Nekov\'a\v{r} und Scholl haben eine ,,plektische'' Erweiterung der Hodge-Theorie entwickelt, die plektische Galois-Aktionen und plektische Hodge-Strukturen einführt. Eine \textit{plektische Gross-Zagier-Formel} würde $L^{(r)}(E,1)$ durch den plektischen Regulator ausdrücken und direkt Hypothese (H1) herstellen.
\end{remark}

\begin{remark}[$p$-adische Deformation nicht-diagonaler Heegner-Zykel]
\label{rem:p-adic-deformation}
Darmons Programm der Stark-Heegner-Punkte konstruiert $p$-adische Punkte auf $E$ aus $p$-adischer Integration auf Shimura-Kurven. Für Rang~$\geq 2$ könnte man Familien nicht-diagonaler Zykel auf höherdimensionalen Shimura-Varietäten $p$-adisch deformieren. Die zentrale Obstruktion ist die Rationalität: ein Abstiegsargument wird benötigt.
\end{remark}

\subsubsection{Euler-Systeme für höheren Rang}

\begin{remark}[Beilinson-Flach-Elemente und Rankin-Selberg-Faltungen]
\label{rem:beilinson-flach}
Kings, Loeffler und Zerbes~\cite{KLZ2017} konstruieren Euler-Systeme aus Beilinson-Flach-Elementen für Rankin-Selberg-Faltungen. Für Rang-2-Kurven könnte man ein ,,Rang-2-Kolyvagin-System'' aus der Kombination von Katos Klassen und Beilinson-Flach-Elementen konstruieren.
\end{remark}

\begin{remark}[Höhere Kolyvagin-Systeme und Selmer-Aufspaltung]
\label{rem:kolyvagin-higher}
Mazur und Rubin~\cite{MazurRubin2004} entwickelten die abstrakte Theorie der Kolyvagin-Systeme. Burns und Sano~\cite{BurnsSano2021} erweiterten diese auf höherrangige Systeme über allgemeinen Koeffizientenringen, und Burns, Sakamoto und Sano~\cite{BurnsSakamotoSano2025} bewiesen die Existenz einer kanonischen höheren Kolyvagin-Ableitungsabbildung zwischen Euler- und Kolyvagin-Systemen in beliebigem Rang, wie von Mazur und Rubin vermutet. Für Rang~$r$ benötigt man ein ,,Rang-$r$-Kolyvagin-System''. Für CM-Kurven liefert der Endomorphismenring eine natürliche Aufspaltung der Selmer-Gruppe in Eigenräume.
\end{remark}

\begin{remark}[Jüngster Fortschritt für Rang $\geq 2$]\label{rem:recent-bsd}
Bedeutender Fortschritt in Richtung höherrangiger BSD wurde in den letzten Jahren erzielt. Burungale, Castella und Kim~\cite{BurungaleCastellaKim2021} bewiesen die Heegner-Punkt-Hauptvermutung von Perrin-Riou, aufbauend auf Howards bipartiten Euler-Systemen und Wei Zhangs Arbeit zu Kolyvagins Vermutung, und etablierten die Iwasawa-Hauptvermutung für die Tate-Shafarevich-Gruppe über antizyklotomischen Erweiterungen. Kim~\cite{Kim2024} erhielt eine höhere Gross-Zagier-Formel durch Vergleich von Heegner-Punkt-Kolyvagin-Systemen mit Kurihara-Zahlen und verfeinerte das strukturelle Verständnis von Selmer-Gruppen. 2024 konstruierten Loeffler und Zerbes \cite{LoefflerZerbes2024} ein universelles Euler-System für $\mathrm{GSp}(4)$ und zeigten, dass alle Euler-System-Klassen in einem eindimensionalen Raum liegen. Burungale, Skinner, Tian und Wan \cite{BSTW2024} bewiesen die Existenz $p$-adischer Zeta-Elemente für elliptische Kurven über imaginär-quadratischen Zahlkörpern, etablierten die Kobayashi-Vermutung für semistabile Kurven an supersingularen Primstellen und präsentieren die ersten unendlichen Familien nicht-CM-elliptischer Kurven, die die volle BSD-Vermutung erfüllen.
\end{remark}

\subsubsection*{Schwellenaxiom: höheres Gross--Zagier}

\begin{axiom}[HGZ-Brücke -- HGZ-BSD]\label{ax:hgz-bsd}
Für eine elliptische Kurve $E/\mathbb{Q}$ vom analytischen Rang $r = \mathrm{ord}_{s=1} L(E,s)$ existiert:
\begin{enumerate}
\item[(HGZ-A)] \textbf{Kanonische Zykel:} Eine Familie arithmetischer Zykel $\mathcal{Z}_r^{(i)}$ ($i = 1, \ldots, r$) in der entsprechenden Chow-Gruppe oder motivischen Kohomologie, funktoriell unter Hecke-Operatoren und Basiswechsel.
\item[(HGZ-B)] \textbf{Höhenkompatibilität:} Eine kanonische Höhenpaarung $\langle \cdot, \cdot \rangle_{\mathrm{ht}}$ auf diesen Zykeln, kompatibel mit der N\'eron--Tate-Höhe für $r=1$ und mit $\det(\langle \mathcal{Z}_r^{(i)}, \mathcal{Z}_r^{(j)} \rangle_{\mathrm{ht}}) = R_E \cdot (\text{lokale Faktoren})$.
\item[(HGZ-C)] \textbf{Selmer-Kontrolle:} Ein Euler-System vom Rang $r$, das die Selmer-Gruppe $\mathrm{Sel}(E/\mathbb{Q})$ kontrolliert und $|\mathrm{Sha}(E/\mathbb{Q})| < \infty$ sowie die Selmer-seitige Ungleichung $\mathrm{rank} \leq \mathrm{ord}$ liefert. Die komplementäre Ungleichung $\mathrm{ord} \leq \mathrm{rank}$ muss aus dem nichtdegenerierten Höhen-/Zykelteil (HGZ-A/B) kommen, nicht aus Selmer-Kontrolle allein.
\end{enumerate}
Die Konjunktion (HGZ-A)+(HGZ-B)+(HGZ-C) impliziert die volle BSD-Formel für Rang $r$.
\end{axiom}

\begin{remark}[Modulstatus]\label{rem:hgz-modules}
\begin{center}
\footnotesize
\resizebox{\columnwidth}{!}{%
\begin{tabular}{@{}lll@{}}
\toprule
\textbf{Modul} & \textbf{Rang 1} & \textbf{Rang $\geq 2$} \\
\midrule
(A) Zykel & Heegner-Punkte (GZ 1986) & Diagonalzykel, partiell \\
(B) Höhen & GZ-Formel (1986) & Offen: höhere GZ-Formel \\
(C) Selmer & Kolyvagin (1988) & Partiell: GSp(4)-Euler-Systeme \\
\bottomrule
\end{tabular}
}
\end{center}
Der kritische Block für Rang $\geq 2$ ist primär \textbf{Modul C}: Existierende Euler-Systeme (Loeffler--Zerbes für GSp(4), Burungale--Skinner--Tian--Wan \cite{BSTW2024} für imaginär-quadratische Zahlkörper) liefern partielle Kontrolle, ergeben aber noch keine vollständige Kolyvagin-artige Schranke für Sha einzelner Kurven vom Rang $\geq 2$.
\end{remark}

\begin{remark}[No-Go-Mechanismen]\label{rem:hgz-nogo}
HGZ-BSD scheitert, wenn:
\begin{enumerate}
\item \textbf{Nicht-kanonische Zykelwahl:} Die höheren Zykel $\mathcal{Z}_r^{(i)}$ hängen von Hilfsdaten (Quaternionenalgebren, Testvektoren) ab, die sich nicht uniformisieren lassen -- dies ist ein Sprachproblem, kein Rechenproblem.
\item \textbf{Degenerierte Höhen:} Die Höhenpaarung $\langle \cdot, \cdot \rangle_{\mathrm{ht}}$ hat einen nicht-trivialen Kern für $r \geq 2$, was die Determinante kollabieren lässt. Höhensättigung (Vermutung~\ref{conj:height_saturation}) schützt dagegen.
\item \textbf{Euler-System-Schwäche:} Das Rang-$r$-Euler-System hat unzureichende Variabilität, um alle $p$-primären Komponenten von Sha gleichzeitig zu kontrollieren.
\end{enumerate}
\end{remark}

\begin{remark}[Modul C als arithmetische Unterbestimmtheit]\label{rem:module-c-underdetermined}
Ohne Rang-$r$-Euler-System-Kontrolle (Modul~C) bleibt die Höhere Gross--Zagier-Formel \emph{arithmetisch unterbestimmt}: Höhen können groß sein, ohne Selmer-Struktur zu implizieren. Dies ist der präzise Inhalt des roten Blocks: Das Brückenlemma (HGZ-A)+(HGZ-B)+(HGZ-C) $\Rightarrow$ $L^{(r)} \sim \det(\text{Höhen})$ erfordert alle drei Module, und Modul~C ist das einzige ohne selbst ein Teilresultat für allgemeine Kurven vom Rang $r \geq 2$. Loeffler--Zerbes \cite{LoefflerZerbes2024} liefern partielle GSp(4)-Kontrolle, aber ein vollständiges Kolyvagin-Ableitungsargument für $|\mathrm{Sha}| < \infty$ bei Rang $\geq 2$ bleibt offen.
\end{remark}

\subsubsection{Arakelov- und motivische Perspektiven}

\begin{remark}[Universeller Arakelov-Regulator]
\label{rem:arakelov-regulator}
Der Arakelov-Regulator $\hat{R}_E$ vereinigt die archimedische N\'eron-Tate-Höhe mit $p$-adischen lokalen Beiträgen zu einer adelischen Höhe. Ein ,,universeller Arakelov-Regulator'' als Dach über $p$-adischem und archimedischem Regulator würde die natürliche Brücke für Hypothese (H1) liefern.
\end{remark}

\begin{remark}[Arakelov-motivische Volumenformel]
\label{rem:arakelov-volume}
Man kann versuchen, $L^{(r)}(E,1)$ als Mass eines ,,arithmetischen Volumens'' eines Arakelov-Vektorbündels auf der arithmetischen Fläche $\mathcal{E}$ zu interpretieren. Die BSD-Formel besagt dann, dass die arithmetische Euler-Charakteristik dem motivischen Volumen entspricht.
\end{remark}

\begin{remark}[Hodge-Arakelov-Volumina: $\Sha$ als topologisches Volumen]
\label{rem:hodge-arakelov-sha}
In Mochizukis Hodge-Arakelov-Theorie könnte $|\Sha|$ als Mass eines topologischen Volumens betrachtet werden, dessen Endlichkeit durch Positivitäts-Constraints (analog den Hodge-Riemann-Relationen) erzwungen wird.
\end{remark}

\subsubsection{Algebraische und kategoriale Ansätze}

\begin{remark}[Höhere Chow-Gruppen auf Picard-Modulflächen]
\label{rem:chow-groups}
Die derivierte Schnittpaarung auf höheren Chow-Gruppen $\mathrm{CH}^r(X, r)$ auf Picard-Modulflächen verallgemeinert die N\'eron-Tate-Paarung und könnte die $r \times r$-Höhenmatrix für Hypothese~(H1) liefern. Die Verbindung zu $L$-Funktionen läuft über die Beilinson-Vermutungen.
\end{remark}

\begin{remark}[Derivierte Hecke-Algebren und Taylor-Wiles-Defekt]
\label{rem:derived-hecke}
Venkatesh' Programm derivierter Hecke-Algebren verbindet die Taylor-Wiles-Methode mit derivierter algebraischer Geometrie. Der ,,Taylor-Wiles-Defekt'' wird vermutet, dem Rang der Mordell-Weil-Gruppe zu entsprechen, was einen neuen Zugang zu BSD eröffnen würde.
\end{remark}

\subsubsection{Analytische und strukturelle Ansätze}

\begin{remark}[$L$-Funktions-Zerlegung in Galois-Sektoren]
\label{rem:galois-sectors}
Die $L$-Funktion zerlegt sich unter Twists durch Charaktere: $L(E, \chi, s)$. Man könnte das BSD-Problem sektorweise lösen (für jeden Twist einzeln) und dann zusammensetzen. Darmon und Rotger~\cite{DarmonRotger2017} verfolgen diese Strategie für Artin-Darstellungen.
\end{remark}

\begin{remark}[Kohomologische Sweeping-Prozesse für $\Sha$-Endlichkeit]
\label{rem:sweeping-sha}
Ein ,,Sweeping-Prozess'' aus der konvexen Analysis würde versuchen zu zeigen, dass der Kern der Global-zu-Lokal-Abbildung in der Galois-Kohomologie endlich ist, indem man iterativ lokale Obstruktionen ,,durchfegt'': an jeder Primstelle $v$ absorbiert der lokale Beitrag $H^1(\mathbb{Q}_v, E)$ einen bestimmten Anteil der globalen Klassen.
\end{remark}

\begin{remark}[O-minimale Strukturen und $\Sha$-Endlichkeit]
\label{rem:o-minimal-sha}
Die O-Minimalmethoden von Bakker-Brunebarbe-Tsimerman~\cite{BBT2023} (erfolgreich auf die Hodge-Vermutung angewendet, vgl.\ das Begleitpaper zu Hodge) könnten potenziell auf $\Sha$-Endlichkeit angewendet werden. Die Zahmheitseigenschaften o-minimaler Strukturen (endliche Stratifikation, kontrollierte Topologie) könnten den Lokus unendlicher $\Sha$ einschränken und möglicherweise als leer zeigen.
\end{remark}

\subsection{Numerische Verifikation}
\label{sec:numerical}

Die BSD-Formel wurde für Tausende elliptischer Kurven bis Rang~4 zu hoher Präzision verifiziert (LMFDB~\cite{LMFDB}). Die folgenden Rechnungen sind bedingte Diagnostik: Sie prüfen Konsistenz mit BSD und mit dem Positivitäts-Normalform-Bild, beweisen aber weder (H1), noch (H2), noch (H3). Unter BSD ist zu erwarten, dass der normalisierte Leitterm den Faktor $|\Sha|/|E_{\mathrm{tors}}|^2$ trifft; der Regulator $R_E$ liefert dabei nur eine bedingte Konsistenzskala, sobald die übrigen BSD-Faktoren festliegen. Die begleitenden Beispiele~\ref{ex:11a}--\ref{ex:389a} verifizieren diese Diagnostik für spezifische Kurven der Ränge~0, 1 und~2.

\subsection{Post-Zookeeper-Transfer: ein Höher-Rang-Modulprogramm}
\label{subsec:zookeeper-bsd-transfer}

Die Post-Zookeeper-Rolle dieses Papers ist nicht, einen neuen Beweis der
BSD-Vermutung in höherem Rang zu behaupten.  Sie besteht darin, die
Höher-Rang-Obstruktion in ein modulares Transferproblem zu zerlegen.  Die
fünf relevanten Felder sind:

Insbesondere hat die CORE-Schliessung von Riemann-$C2/MS2$ keine
automatische BSD-Konsequenz.  Sie fixiert nur die gemeinsame
Normalform-Disziplin: Operator/Form, expliziter Defekt, endliche
Approximation, Gap und Grenzübergang müssen getrennt ausgewiesen werden.
Für BSD bleibt der rote Block die höher-rangige Selmer-Kontrolle und die
Höhere-Gross--Zagier-Brücke, insbesondere Modul~C in
Bemerkung~\ref{rem:hgz-modules}.

\begin{description}
  \item[Operator/Form.] Die Neron--Tate- und $p$-adischen Höhenpaarungen
  auf Mordell--Weil- und Selmer-Richtungen.
  \item[Defekt.] Der BSD-Defekt
  \[
    \delta_{\mathrm{BSD}}(E)
      := \operatorname{ord}_{s=1}L(E,s)-\operatorname{rank}E(\mathbb{Q}),
  \]
  zusammen mit dem Residuum zwischen vorhergesagtem Leitterm und
  Regulatorausdruck.
  \item[Approximation.] Endlicher Leiter, fixer Rang, endliches
  Selmer-Level und endliche Stufen eines Iwasawa-Turms.
  \item[Gap.] Nichtausartung des archimedischen oder $p$-adischen
  Regulators.
  \item[Grenzübergang.] Der Übergang von $p$-adischer
  Leitterminformation und Selmer-Struktur zur archimedischen
  BSD-Leittermformel.
\end{description}

Daraus ergibt sich ein konservatives nächstes Ziel.  Statt einen
allgemeinen Rang-$r$-Beweis zu formulieren, sollte zuerst ein
Rang-zwei-Transfersatz isoliert werden:
\[
\begin{gathered}
  \text{höherer Gross--Zagier-Input}\\
  +\ \text{Selmer-Kontrolle}\\
  +\ \text{$p$-adische Höhen-Nichtausartung}\\
  \Longrightarrow \text{Rang-zwei-Höhensättigung}.
\end{gathered}
\]
Kims Arbeiten zu höherem Gross--Zagier und Selmer/Iwasawa sowie der
$p$-adische BSD-Rahmen von Castella et al. liefern dafür den natürlichen
Literatur-Backbone.  Der erwartete Gewinn ist diagnostisch, nicht final:
Höherer Rang wird in vier separat testbare Module zerlegt, statt eine
einzige monolithische Lücke zu bleiben.

\begin{remark}[Determinantenlinien-Diagnostik und Advice-Guardrail]
\label{rem:det-line-diagnostic}
Das aktuelle Arbeitsledger schärft das Rang-zwei-Programm, indem es den
Regulator als Norm auf einer Determinantenlinie und nicht als nachträglichen
Skalar behandelt. Für ein Kandidatenpaket $Z=(Z_1,\ldots,Z_r)$ setze
\[
H_Z=(\langle Z_i,Z_j\rangle_{\mathrm{NT}})_{i,j},\qquad
G_r(E,Z)=\frac{\det H_Z}{\prod_i (H_Z)_{ii}},
\]
ergänzt durch den Degenerationsindikator
$\widetilde G_r(E,Z)=\lambda_{\min}(H_Z)/\mathrm{tr}(H_Z)$, sofern die
Einträge verfügbar sind. Das vorgeschlagene Defektledger ist
\[
\begin{aligned}
D_{\det}(E,Z)=(&d_{\mathrm{analytic\ line}},d_{\mathrm{compare}},d_{\mathrm{gap}},\\
&d_{\mathrm{selmer\ tail}},d_{\mathrm{advice}}).
\end{aligned}
\]
Die letzte Komponente hält fest, ob die Kanalwahl aus analytischen,
$p$-adischen oder Hecke-Daten stammt oder ob sie eine bekannte
Mordell--Weil-Basis bzw. manuelle Ranginformation benutzt. Letzteres ist als
Referenznormalisierung nützlich, darf aber nicht als Brückenevidenz zählen.
Leitterm-, Regulator- und Mordell--Weil-Information dürfen also nicht
zugleich Zielgröße und Prädiktor sein; sonst recycelt die Diagnostik genau
die BSD-Größe, die sie testen soll.

Ein Rang-zwei-Signal wird erst gezählt, wenn dieselbe Pipeline an
vorab registrierten Negativkontrollen scheitert: duplizierte Kanäle
$Z_2=Z_1$, Rang-eins-Kurven als Rang-zwei-Kandidaten eingebettet und
führerabgestimmte Kontrollen mit Kanalregeln, die vor dem analytischen
Vergleich festgelegt wurden.
\end{remark}

\subsection*{Ergebnisklassifikation}

Dieses Paper etabliert eine \textbf{strukturelle Normalform-Charakterisierung innerhalb der BSD/Bloch--Kato-motivischen Produktklasse}: unter vier Positivitätsaxiomen plus den expliziten Brückenhypothesen (H1)--(H3) ist die BSD-Formel die erwartete zulässige kanonische Produktgestalt für den Leitterm von $L^{(r)}(E,1)$. Es ist kein Eindeutigkeitsclaim unter beliebigen Funktionen. Für Rang~$\leq 1$ ist dies eine genuine Umformulierung bewiesener Resultate; für Rang~$\geq 2$ ist es eine bedingte strukturelle Analyse.

\begin{center}
\footnotesize
\resizebox{\columnwidth}{!}{%
\begin{tabular}{@{}ll@{}}
\toprule
\textbf{Ergebnis} & \textbf{Typ} \\
\midrule
\multicolumn{2}{@{}l}{\textit{Bewiesen (unbedingt):}} \\
\midrule
N\'eron--Tate-Positivität (Axiom II) & Klassischer Satz \\
BSD-Konsequenzen für Rang $\leq 1$ & Satz; Normalform-Lesart interpretativ \\
\midrule
\multicolumn{2}{@{}l}{\textit{Numerische Evidenz (kein Beweis):}} \\
\midrule
Rang-2-Konsistenzdiagnostik & Bedingte Checks; Advice-Audit offen \\
Regulator-Skalierung $R \sim N^{0.10}$ & Explorative Regression \\
\midrule
\multicolumn{2}{@{}l}{\textit{Strukturelle Beiträge:}} \\
\midrule
Axiomensystem (I--IV) & Definition / strukturelle Bedingungen \\
Eindeutigkeit der BSD-Produktgestalt & Produktklassen-Claim; bedingt durch H1--H3 \\
Gross--Zagier als Prototyp & Interpretation \\
\midrule
\multicolumn{2}{@{}l}{\textit{Offen (fundamental):}} \\
\midrule
Höheres Gross--Zagier ($r \geq 2$) & Offen \\
Höhere Kolyvagin-Systeme & Offen \\
Sha-Endlichkeit ($r \geq 2$) & Offen \\
\bottomrule
\end{tabular}
}
\end{center}

% ===================================================================
% 11. SCHLUSSFOLGERUNG
% ===================================================================
\section{Schlussfolgerung}
\label{sec:conclusion}

Die Birch-und-Swinnerton-Dyer-Vermutung kann in dieser strukturellen Lesart als Diskrepanzkontrollproblem zwischen analytischen und geometrischen Strukturen verstanden werden. Die N\'eron-Tate-Höhe liefert die kanonische Positivitätsstruktur: sie ist eine positiv-definite Form auf rationalen Punkten, die arithmetische Richtungen über ihre Energiekosten detektiert.

Die BSD-Formel ist die erwartete kanonische Produktnormalform, die mit vier axiomatischen Zwangsbedingungen plus den nötigen Brückenhypothesen kompatibel ist: endliche arithmetische Kapazität ($\Sha$-Endlichkeit), kooperative Höhenpositivität (N\'eron-Tate), kontrolliertes Iwasawa-Wachstum, Euler-System-Kompositionskohärenz, Höhenkopplung und Selmer-Kontrolle. Das reduziert BSD nicht auf Positivität allein; in höherem Rang werden BSD-starke Eingaben in benannte Brückenmodule ausgelagert.

Der Rang-$\leq 1$-Beweis ist mit dieser Positivitäts-No-Go-Lesart kompatibel: Gross-Zagier koppelt eine analytische Ableitung an eine nichtnegative Höhe, und Kolyvagins Euler-System liefert Selmer-Kontrolle. Für höheren Rang identifiziert das Framework zwei fehlende Zutaten: nichtdegenerierte höhere Gross-Zagier-Höhen-/Zykelkopplung und volle höher-rangige Selmer-/Kolyvagin-/Sha-Kontrolle. Height Saturation benennt das Ziel, schließt aber keine der beiden Brücken.

Für Rang $\leq 1$ ist der Beitrag konzeptionell: eine Normalform-Lesart bekannter Theoreme. Für Rang $\geq 2$ ist der Beitrag ein abgesichertes Forschungsprogramm und ein Diagnoseledger, das Theorem-Eingaben, konjekturale Brückeneingaben und advice-normalisierte numerische Checks trennt.

\begin{quote}
\textit{,,BSD verlangt nicht von uns, rationale Punkte zu finden. Es verlangt von uns zu zeigen, dass die $L$-Funktion keinen Raum für Nullstellen ohne sie hat.''}
\end{quote}

\section*{KI-Offenlegung}
KI-basierte Werkzeuge wurden zur Unterstützung bei mathematischer Formalisierung, Literaturverifikation, Übersetzung, Konsistenzprüfung und redaktioneller Überarbeitung eingesetzt. Sie waren keine Autorinnen oder Autoren, keine Koautoren, keine Danksagungsempfänger, keine Ground-Truth-/Wahrheitsinstanzen und keine unabhängigen Validierungsinstanzen. Forschungsrichtung, mathematische Aussagen, finale redaktionelle Entscheidungen und Verantwortung für die Korrektheit liegen ausschließlich beim Autor.


% ===================================================================
% DANKSAGUNGEN
% ===================================================================
\begin{acknowledgments}
Diese Arbeit wurde als unabhängige Forschung ohne institutionelle Förderung durchgeführt.

\textit{Förderinformation.} Diese Forschung erhielt keine externe Förderung.
\end{acknowledgments}


% ===================================================================
% BIBLIOGRAPHIE
% ===================================================================
\begin{thebibliography}{35}
\setlength{\itemsep}{0pt}
\setlength{\parskip}{0pt}

\bibitem{BSD1965}
B.~J.~Birch \& H.~P.~F.~Swinnerton-Dyer (1965).
Notes on elliptic curves. II.
\textit{Journal f\"ur die reine und angewandte Mathematik}, 218, 79--108.
DOI: 10.1515/crll.1965.218.79.

\bibitem{GrossZagier1986}
B.~Gross \& D.~Zagier (1986).
Heegner points and derivatives of $L$-series.
\textit{Inventiones Mathematicae}, 84, 225--320.
DOI: 10.1007/BF01388809.

\bibitem{Kolyvagin1990}
V.~A.~Kolyvagin (1990).
Euler systems.
In \textit{The Grothendieck Festschrift}, Vol.~II, pp.\ 435--483. Birkh\"auser.
DOI: 10.1007/978-0-8176-4575-5\_11.

\bibitem{Kato2004}
K.~Kato (2004).
$p$-adic Hodge theory and values of zeta functions of modular forms.
\textit{Ast\'erisque}, 295, 117--290.

\bibitem{Kim2024}
C.-H.~Kim (2024).
A higher Gross-Zagier formula and the structure of Selmer groups.
Erscheint in \textit{Transactions of the AMS}.
arXiv:2203.12161.

\bibitem{SkinnerUrban2014}
C.~Skinner \& E.~Urban (2014).
The Iwasawa main conjectures for $\mathrm{GL}_2$.
\textit{Inventiones Mathematicae}, 195(1), 1--277.
DOI: 10.1007/s00222-013-0448-1.

\bibitem{MazurRubin2004}
B.~Mazur \& K.~Rubin (2004).
\textit{Kolyvagin Systems}.
Memoirs of the AMS, 168(799).
DOI: 10.1090/memo/0799.

\bibitem{DarmonRotger2017}
H.~Darmon \& V.~Rotger (2017).
Diagonal cycles and Euler systems II: The Birch and Swinnerton-Dyer conjecture for Hasse-Weil-Artin $L$-functions.
\textit{Journal of the AMS}, 30(3), 601--672.
DOI: 10.1090/jams/861.

\bibitem{BDP2013}
M.~Bertolini, H.~Darmon, \& K.~Prasanna (2013).
Generalized Heegner cycles and $p$-adic Rankin $L$-series.
\textit{Duke Mathematical Journal}, 162(6), 1033--1148.
DOI: 10.1215/00127094-2142056.

\bibitem{BurungaleCastellaKim2021}
A.~Burungale, F.~Castella \& C.-H.~Kim (2021).
A proof of Perrin-Riou's Heegner point main conjecture.
\textit{Algebra \& Number Theory}, 15(7), 1627--1653.
DOI: 10.2140/ant.2021.15.1627.

\bibitem{BurnsSano2021}
D.~Burns \& T.~Sano (2021).
On the theory of higher rank Euler, Kolyvagin and Stark systems.
\textit{International Mathematics Research Notices}, 2021(13), 10118--10206.
DOI: 10.1093/imrn/rnz103.

\bibitem{BurnsSakamotoSano2025}
D.~Burns, R.~Sakamoto \& T.~Sano (2025).
On the theory of higher rank Euler, Kolyvagin and Stark systems, II: the general theory.
Erscheint in \textit{Algebra \& Number Theory}.
arXiv:1805.08448.

\bibitem{LLZ2014}
A.~Lei, D.~Loeffler, \& S.~L.~Zerbes (2014).
Euler systems for Rankin-Selberg convolutions of modular forms.
\textit{Annals of Mathematics}, 180(2), 653--771.
DOI: 10.4007/annals.2014.180.2.6.

\bibitem{LSZ2020}
D.~Loeffler, C.~Skinner, \& S.~L.~Zerbes (2022).
Euler systems for $\mathrm{GSp}(4)$.
\textit{Journal of the European Mathematical Society}, 24(2), 669--733.
DOI: 10.4171/JEMS/1124.

\bibitem{Silverman2009}
J.~H.~Silverman (2009).
\textit{The Arithmetic of Elliptic Curves}, 2nd ed.
Graduate Texts in Mathematics 106, Springer.
DOI: 10.1007/978-0-387-09494-6.

\bibitem{Wiles1995}
A.~Wiles (1995).
Modular elliptic curves and Fermat's Last Theorem.
\textit{Annals of Mathematics}, 141(3), 443--551.
DOI: 10.2307/2118559.

\bibitem{Geiger2026e}
L.~Geiger (2026).
The Saturation Theorem: Axiomatic Constraints on Quantum Gravity from Cosmological Relaxation.
Zenodo preprint.
\href{https://doi.org/10.5281/zenodo.19036188}{doi:10.5281/zenodo.19036188}.

\bibitem{FST-Unified2026}
L.~Geiger (2026).
FST Spectrum Duality: The Renormalized Free Energy Principle as Physical Instantiation of Functional Stability Theory.
Zenodo preprint.
\href{https://doi.org/10.5281/zenodo.19036190}{doi:10.5281/zenodo.19036190}.

\bibitem{GeigerRH2026c}
L.~Geiger (2026).
From Landscape to Proof: The Even-Dominance Route to the Riemann Hypothesis.
Zenodo preprint.
\href{https://doi.org/10.5281/zenodo.19035640}{doi:10.5281/zenodo.19035640}.

\bibitem{GreenbergStevens1993}
R.~Greenberg \& G.~Stevens (1993).
$p$-adic $L$-functions and $p$-adic periods of modular forms.
\textit{Inventiones Mathematicae}, 111(2), 407--447.
DOI: 10.1007/BF01231294.

\bibitem{BhargavaShankar2015}
M.~Bhargava \& A.~Shankar (2015).
Ternary cubic forms having bounded invariants, and the existence of a positive proportion of elliptic curves having rank~0.
\textit{Annals of Mathematics}, 181(2), 587--621.
DOI: 10.4007/annals.2015.181.2.4.

\bibitem{LMFDB}
The LMFDB Collaboration (2024).
\textit{The L-functions and Modular Forms Database}.
\url{https://www.lmfdb.org}.

\bibitem{Nekovar2006}
J.~Nekov\'a\v{r} (2006).
Selmer complexes.
\textit{Ast\'erisque}, 310, vi+559.

\bibitem{Dokchitser2010}
T.~Dokchitser \& V.~Dokchitser (2010).
On the Birch-Swinnerton-Dyer quotients modulo squares.
\textit{Annals of Mathematics}, 172(1), 567--596.
DOI: 10.4007/annals.2010.172.567.

\bibitem{YZZ2013}
X.~Yuan, S.-W.~Zhang, \& W.~Zhang (2013).
\textit{The Gross-Zagier Formula on Shimura Curves}.
Annals of Mathematics Studies 184, Princeton University Press.

\bibitem{PerrinRiou1987}
B.~Perrin-Riou (1987).
Points de Heegner et d\'eriv\'ees de fonctions $L$ $p$-adiques.
\textit{Inventiones Mathematicae}, 89(3), 455--510.
DOI: 10.1007/BF01388982.

\bibitem{KLZ2017}
G.~Kings, D.~Loeffler, \& S.~L.~Zerbes (2017).
Rankin-Eisenstein classes and explicit reciprocity laws.
\textit{Cambridge Journal of Mathematics}, 5(1), 1--122.
DOI: 10.4310/CJM.2017.v5.n1.a1.

\bibitem{BBT2023}
B.~Bakker, B.~Brunebarbe \& J.~Tsimerman (2023).
o-minimal GAGA and a conjecture of Griffiths.
\textit{Inventiones Mathematicae}, 232(1), 163--228.
DOI: 10.1007/s00222-022-01166-1.

\bibitem{Smith2022}
A.~Smith, \textit{$2^\infty$-Selmer-Gruppen, $2^\infty$-Klassengruppen und Goldfelds Vermutung}, arXiv:1702.02325v3, 2022.

\bibitem{LoefflerZerbes2024} D.~Loeffler und S.~L.~Zerbes, \textit{A universal Euler system for GSp(4)}, arXiv:2411.12576, 2024.

\bibitem{BSTW2024} A.~Burungale, C.~Skinner, Y.~Tian und X.~Wan, \textit{Zeta elements for elliptic curves and applications}, arXiv:2409.01350, 2024.

\bibitem{Cesnavicius2018}
K.~\v{C}esnavi\v{c}ius (2018).
The Manin constant in the semistable case.
\textit{Compositio Mathematica}, 154(9), 1889--1920.
DOI: 10.1112/S0010437X18007273.

\bibitem{BlochKato1990}
S.~Bloch \& K.~Kato (1990).
$L$-functions and Tamagawa numbers of motives.
In \textit{The Grothendieck Festschrift}, Vol.~I, pp.\ 333--400. Birkh\"auser.
DOI: 10.1007/978-0-8176-4574-8\_9.

\bibitem{FontainePerrinRiou1994}
J.-M.~Fontaine \& B.~Perrin-Riou (1994), \textit{Autour des conjectures de Bloch et Kato: cohomologie galoisienne et valeurs de fonctions~$L$}, in \textit{Motives} (Seattle 1991), Proc.\ Symp.\ Pure Math.\ 55, Part~1, 599--706. AMS.

\bibitem{Szpiro1981}
L.~Szpiro (1981), \textit{Propri\'et\'es num\'eriques du faisceau dualisant r\'elatif}, Ast\'erisque \textbf{86}, 44--78.

\bibitem{GeigerIUT2026}
L.~Geiger (2026).
\textit{The Gap in IUTchIII, Corollary~3.12: From Attempted Repair to Structural Diagnosis (DRAFT)}.
Zenodo preprint.
\href{https://doi.org/10.5281/zenodo.19960781}{doi:10.5281/zenodo.19960781}.

\end{thebibliography}

\end{document}


